%%%%%%%%%%%%%%%%%%%%%%%%%%%%%%%%%%%%%%%%%%%%%%%%%%%%%%%%%%%%%%%%%%%%%%%%%%%%%%%
%  SÍLABO DE MACROECONOMETRÍA — Plan de estudio integrado
%
%  Textos base:
%    MHH — Martin, Hurn y Harris, Econometric Modelling with Time Series
%          (Cambridge, 2013)
%    BE  — Bårdsen, Eitrheim, Jansen y Nymoen, The Econometrics of
%          Macroeconomic Modelling (Oxford, 2005)
%  Textos de apoyo:
%    En  — Enders, Applied Econometric Time Series (Wiley, 4.ª ed., 2014)
%    Hm  — Hamilton, Time Series Analysis (Princeton, 1994)
%    KL  — Kilian y Lütkepohl, Structural Vector Autoregressive Analysis
%          (Cambridge, 2017)
%    BM  — Blake y Mumtaz, Applied Bayesian Econometrics for Central
%          Bankers (Bank of England, ed. actualizada 2017)
%    Fa  — Fair, Macroeconomic Modeling: The Cowles Commission Approach
%          (MIT Press, 2024)
%    ET  — Elliott y Timmermann, Economic Forecasting (Princeton, 2016)
%          (lectura principal del Bloque 15, de pronóstico)
%    Ps  — Pesaran, Time Series and Panel Data Econometrics (Oxford, 2015)
%  Programas PUCP integrados: 1ECO24 (2025-1, 2025-2, 2026-1),
%    1ECO78 (2024-2, 2025-2, 2026-2), 2ECO01 (unidad 4); además, el
%    programa de Macroeconometría del Curso de Extensión del BCRP (2004).
%
%  Compilar con:  pdflatex Macroeconometria.tex   (dos veces, por el índice)
%%%%%%%%%%%%%%%%%%%%%%%%%%%%%%%%%%%%%%%%%%%%%%%%%%%%%%%%%%%%%%%%%%%%%%%%%%%%%%%

\documentclass[11pt,a4paper]{article}

\usepackage[T1]{fontenc}
\usepackage[utf8]{inputenc}
\usepackage[spanish,es-noquoting,es-tabla]{babel}
\usepackage{lmodern}
\usepackage[scaled=0.92]{helvet}        % títulos en sans serif, al estilo McGraw-Hill
\usepackage[margin=2.4cm,headheight=15pt]{geometry}
\usepackage[table]{xcolor}
\usepackage{enumitem}
\usepackage{titlesec}
\usepackage{titletoc}
\usepackage{tabularx}
\usepackage{booktabs}
\usepackage{longtable}
\usepackage{amsmath}
\usepackage{amssymb}
\usepackage{fancyhdr}
\usepackage{tikz}
\usepackage[skins]{tcolorbox}
\usepackage{needspace}

%------------------------------------------------------------------ colores ---
% Gama de turquesas
\definecolor{tonoOscuro}{RGB}{4,75,82}              % oscuro
\definecolor{tonoProfundo}{RGB}{6,110,121}          % profundo
\definecolor{tonoPrincipal}{RGB}{8,145,160}         % principal
\definecolor{tonoA}{RGB}{25,133,179}                % variante cálida
\definecolor{tonoB}{RGB}{15,179,168}                % variante fría
\definecolor{tonoC}{RGB}{21,193,153}                % variante viva
\definecolor{tonoClaro}{RGB}{95,208,221}            % claro
\definecolor{tonoAcero}{RGB}{63,124,131}            % acero
\definecolor{tonoPizarra}{RGB}{86,129,134}          % pizarra
\definecolor{tonoGrafito}{RGB}{39,70,73}            % grafito (texto secundario)
\definecolor{tonoHielo}{RGB}{207,243,247}           % hielo
\definecolor{tonoPalido}{RGB}{233,249,251}          % pálido
\definecolor{tonoBruma}{RGB}{246,253,254}           % bruma

\colorlet{azulosc}{tonoOscuro}
\colorlet{grisclaro}{tonoPalido}
% Un tono distinto de la gama para cada libro
\definecolor{colorMHH}{RGB}{15,138,105}
\definecolor{colorBE}{RGB}{20,184,155}
\definecolor{colorEn}{RGB}{18,161,149}
\definecolor{colorHm}{RGB}{15,137,138}
\definecolor{colorKL}{RGB}{20,167,184}
\definecolor{colorBM}{RGB}{18,133,161}
\definecolor{colorFa}{RGB}{15,103,138}
\definecolor{colorET}{RGB}{20,122,184}
\definecolor{colorPs}{RGB}{18,94,161}
\colorlet{colorPUCP}{tonoPizarra}

\usepackage[colorlinks=true,linkcolor=tonoOscuro,urlcolor=tonoPrincipal]{hyperref}

%------------------------------------------------------------------ estilos ---
\titleformat{\section}[block]
  {\sffamily\LARGE\bfseries\color{tonoOscuro}\raggedright}{\thesection.}{0.6em}
  {\textcolor{tonoPrincipal}{\rule[-0.1ex]{0.55em}{0.85em}}\hspace{0.45em}}
  [\vspace{-0.4ex}{\color{tonoPrincipal}\titlerule[1.6pt]}\vspace{1.2pt}{\color{tonoClaro}\titlerule[0.6pt]}]
\titlespacing*{\section}{0pt}{2.6ex plus 1ex minus .2ex}{1.6ex plus .2ex}

\titleformat{\subsection}
  {\sffamily\large\bfseries\color{tonoPrincipal}}{\thesubsection}{0.6em}
  {\textcolor{tonoClaro}{\rule[0.1ex]{0.45em}{0.45em}}\hspace{0.45em}}
\titlespacing*{\subsection}{0pt}{2.2ex plus .8ex minus .2ex}{1ex plus .2ex}

% Índice
\titlecontents{section}[0pt]
  {\addvspace{3pt}\sffamily\bfseries\color{tonoOscuro}}
  {}{\textcolor{tonoB}{\rule[0.15ex]{0.4em}{0.4em}}\hspace{0.6em}}
  {{\color{tonoClaro}\titlerule*[0.7em]{.}}\textcolor{tonoPrincipal}{\contentspage}}

% Cabecera y pie
\pagestyle{fancy}
\fancyhf{}
\lhead{\small\sffamily\bfseries\color{tonoOscuro}Sílabo de Macroeconometría}
\rhead{\footnotesize\sffamily\color{tonoPrincipal}Martin--Hurn--Harris $\cdot$ Bårdsen et al. $\cdot$ Enders $\cdot$ Hamilton $\cdot$ Kilian--Lütkepohl}
\cfoot{{\setlength{\fboxsep}{3pt}\colorbox{tonoPrincipal}{\makebox[1.8em]{\small\sffamily\bfseries\color{white}\thepage}}}}
\renewcommand{\headrule}{%
  {\color{tonoPrincipal}\hrule height 1.1pt width\headwidth}\vskip 1.2pt
  {\color{tonoClaro}\hrule height 0.4pt width\headwidth}\vskip -2.7pt}

% Listas
\setlist[itemize,1]{label=\textcolor{tonoPrincipal}{\scriptsize$\blacksquare$}}
\setlist[enumerate,1]{label=\textcolor{tonoPrincipal}{\sffamily\bfseries\arabic*.}}

% Tablas: cabecera oscura y filas alternas en tono pálido
\arrayrulecolor{tonoPrincipal}
\setlength{\aboverulesep}{0pt}
\setlength{\belowrulesep}{0pt}
\renewcommand{\arraystretch}{1.3}
\newcommand{\cab}[1]{\textcolor{white}{\sffamily\bfseries #1}}
\newcommand{\filasalternas}{\rowcolors{2}{white}{tonoPalido}}

% Encabezados de párrafo en color
\newcommand{\destacado}[1]{{\sffamily\bfseries\color{tonoPrincipal}#1}}

% Franja decorativa con la gama
\newcommand{\franjas}[1][3pt]{%
  \noindent\begin{tikzpicture}
    \foreach \c [count=\i from 0] in {tonoOscuro,tonoProfundo,tonoPrincipal,tonoA,tonoB,tonoC,tonoClaro,tonoHielo}
      \fill[\c] (\i*\linewidth/8,0) rectangle ++(\linewidth/8,#1);
  \end{tikzpicture}\par}

%---------------------------------------------------------------- comandos ----
% Etiqueta de libro: recuadro de color con la sigla
\newcommand{\etiqueta}[2]{%
  {\setlength{\fboxsep}{1.4pt}\colorbox{#1}{\sffamily\bfseries\color{white}\strut#2}}}
% Referencias abreviadas a cada libro
\newcommand{\MHH}[1]{\etiqueta{colorMHH}{MHH}\nobreak\hspace{0.25em}\textcolor{colorMHH}{#1}}
\newcommand{\BE}[1]{\etiqueta{colorBE}{BE}\nobreak\hspace{0.25em}\textcolor{colorBE}{#1}}
\newcommand{\EN}[1]{\etiqueta{colorEn}{En}\nobreak\hspace{0.25em}\textcolor{colorEn}{#1}}
\newcommand{\HM}[1]{\etiqueta{colorHm}{Hm}\nobreak\hspace{0.25em}\textcolor{colorHm}{#1}}
\newcommand{\KL}[1]{\etiqueta{colorKL}{KL}\nobreak\hspace{0.25em}\textcolor{colorKL}{#1}}
\newcommand{\BM}[1]{\etiqueta{colorBM}{BM}\nobreak\hspace{0.25em}\textcolor{colorBM}{#1}}
\newcommand{\FA}[1]{\etiqueta{colorFa}{Fa}\nobreak\hspace{0.25em}\textcolor{colorFa}{#1}}
\newcommand{\ET}[1]{\etiqueta{colorET}{ET}\nobreak\hspace{0.25em}\textcolor{colorET}{#1}}
\newcommand{\PS}[1]{\etiqueta{colorPs}{Ps}\nobreak\hspace{0.25em}\textcolor{colorPs}{#1}}

% Un tema: número en recuadro de color + título
\newcommand{\tema}[1]{\temaaux#1\relax}
\def\temaaux#1 #2\relax{%
  \par\addvspace{1.1em}\noindent
  {\sffamily\bfseries\large\raggedright
   \hangindent=2.9em\hangafter=1
   \makebox[2.9em][l]{{\setlength{\fboxsep}{2.5pt}\colorbox{tonoPrincipal}{\color{white}#1}}}%
   {\color{tonoOscuro}#2}\par}%
  \nobreak\vspace{0.25em}\noindent{\color{tonoClaro}\rule{\linewidth}{0.6pt}}\par\nobreak\vspace{0.35em}}

% Línea de lecturas con marcador y etiqueta de color
\newcommand{\linealectura}[3]{%
  \par\nobreak\noindent\hspace{1.2em}%
  \begin{minipage}{0.92\textwidth}\small
    \textcolor{#1}{\rule[0.1ex]{0.45em}{0.45em}}\hspace{0.45em}%
    {\sffamily\bfseries\color{#1}#2}\enspace#3%
  \end{minipage}\par\vspace{0.3em}}
\newcommand{\lecturas}[1]{\linealectura{tonoPrincipal}{Lecturas:}{#1}}
\newcommand{\apoyo}[1]{\linealectura{tonoA}{Profundización:}{#1}}
\newcommand{\juegos}[1]{\linealectura{tonoB}{Teoría de juegos:}{#1}}
\newcommand{\supl}[1]{\linealectura{tonoC}{Suplemento:}{#1}}
\newtcolorbox{cajanota}{enhanced, sharp corners, boxrule=0pt, frame hidden,
  colback=tonoPalido, borderline west={2.5pt}{0pt}{tonoClaro},
  left=8pt, right=6pt, top=4pt, bottom=4pt, before skip=0pt, after skip=0pt,
  fontupper=\small\color{tonoGrafito}}
\newcommand{\nota}[1]{%
  \par\nobreak\noindent\hspace{1.2em}%
  \begin{minipage}{0.92\textwidth}\begin{cajanota}#1\end{cajanota}\end{minipage}\par\vspace{0.3em}}
\newcommand{\consultalib}[1]{\linealectura{tonoPizarra}{Consulta:}{{\color{tonoGrafito}#1}}}
% Artículos y textos tomados de los programas PUCP
\newcommand{\articulos}[1]{\linealectura{colorPUCP}{Artículos (PUCP):}{#1}}
% Artículos que no proceden de los programas PUCP (entran en las horas)
\newcommand{\articulosB}[1]{\linealectura{tonoB}{Artículos:}{#1}}
\newcommand{\consulta}[1]{\linealectura{colorPUCP}{Consulta (PUCP):}{{\color{tonoGrafito}#1}}}

% Encabezado de bloque: banda oscura, franja de la gama y recuadro de presentación
\newif\ifrecienparte
\newcommand{\bloque}[2]{%
  \ifrecienparte\global\recienpartefalse\else\par\needspace{0.4\textheight}\fi
  \phantomsection\addcontentsline{toc}{section}{#1}%
  \begin{tcolorbox}[enhanced, sharp corners, boxrule=0pt, frame hidden,
      colback=tonoOscuro, borderline west={7pt}{0pt}{tonoClaro},
      left=14pt, right=10pt, top=10pt, bottom=10pt, before skip=0pt, after skip=0pt]
    {\raggedright\sffamily\Large\bfseries\color{white}#1\par}
  \end{tcolorbox}%
  \nointerlineskip\franjas[4pt]
  \vspace{0.8em}
  \begin{tcolorbox}[enhanced, sharp corners, boxrule=0pt, frame hidden,
      colback=tonoPalido, borderline west={3pt}{0pt}{tonoPrincipal},
      borderline south={0.6pt}{0pt}{tonoClaro},
      left=10pt, right=8pt, top=6pt, bottom=6pt, before skip=0pt, after skip=0pt,
      fontupper=\small\color{tonoGrafito}]
    #2
  \end{tcolorbox}%
  \par\medskip}

%================================================================== documento ==
\begin{document}

\begin{titlepage}
\begin{tikzpicture}[remember picture,overlay]
  % banda superior
  \fill[tonoOscuro] (current page.north west) rectangle ([yshift=-5.4cm]current page.north east);
  \fill[tonoClaro] (current page.north west) rectangle ([xshift=1.1cm,yshift=-5.4cm]current page.north west);
  % franja de la gama bajo la banda
  \foreach \c [count=\i from 0] in {tonoOscuro,tonoProfundo,tonoPrincipal,tonoA,tonoB,tonoC,tonoClaro,tonoHielo}
    \fill[\c] ([xshift=\i*\paperwidth/8,yshift=-5.4cm]current page.north west)
      rectangle ++(\paperwidth/8,-0.35cm);
  % franja inferior
  \foreach \c [count=\i from 0] in {tonoHielo,tonoClaro,tonoC,tonoB,tonoA,tonoPrincipal,tonoProfundo,tonoOscuro}
    \fill[\c] ([xshift=\i*\paperwidth/8]current page.south west)
      rectangle ++(\paperwidth/8,0.6cm);
\end{tikzpicture}%
\centering
\vspace*{0.5cm}
{\sffamily\Huge\bfseries\color{white} Sílabo de Macroeconometría\par}
\vspace{0.4cm}
{\sffamily\Large\color{tonoHielo} Plan de estudio integrado\par}
\vspace{1.9cm}
{\sffamily\large\bfseries\color{tonoPrincipal}\textcolor{tonoClaro}{\rule[0.5ex]{1.5cm}{0.6pt}}\quad Textos base\quad\textcolor{tonoClaro}{\rule[0.5ex]{1.5cm}{0.6pt}}\par}
\vspace{0.3cm}
{\large\etiqueta{colorMHH}{MHH} \quad {\color{tonoOscuro}Vance Martin, Stan Hurn y David Harris}\par}
{\normalsize \emph{\color{colorMHH}Econometric Modelling with Time Series: Specification, Estimation and Testing}\par}
{\small\color{tonoPizarra} Cambridge University Press, 2013 --- 21 capítulos\par}
\vspace{0.12cm}
{\large\etiqueta{colorBE}{BE} \quad {\color{tonoOscuro}Gunnar Bårdsen, Øyvind Eitrheim, Eilev S. Jansen y Ragnar Nymoen}\par}
{\normalsize \emph{\color{colorBE}The Econometrics of Macroeconomic Modelling}\par}
{\small\color{tonoPizarra} Oxford University Press, 2005 --- 11 capítulos y apéndice\par}
\vspace{0.35cm}
{\sffamily\large\bfseries\color{tonoPrincipal}\textcolor{tonoClaro}{\rule[0.5ex]{1.5cm}{0.6pt}}\quad Textos de apoyo\quad\textcolor{tonoClaro}{\rule[0.5ex]{1.5cm}{0.6pt}}\par}
\vspace{0.3cm}
{\large\etiqueta{colorEn}{En} \quad {\color{tonoOscuro}Walter Enders}\par}
{\normalsize \emph{\color{colorEn}Applied Econometric Time Series}\par}
{\small\color{tonoPizarra} Wiley, 4.ª edición, 2014 --- 7 capítulos\par}
\vspace{0.12cm}
{\large\etiqueta{colorHm}{Hm} \quad {\color{tonoOscuro}James D. Hamilton}\par}
{\normalsize \emph{\color{colorHm}Time Series Analysis}\par}
{\small\color{tonoPizarra} Princeton University Press, 1994 --- 22 capítulos\par}
\vspace{0.12cm}
{\large\etiqueta{colorKL}{KL} \quad {\color{tonoOscuro}Lutz Kilian y Helmut Lütkepohl}\par}
{\normalsize \emph{\color{colorKL}Structural Vector Autoregressive Analysis}\par}
{\small\color{tonoPizarra} Cambridge University Press, 2017 --- 19 capítulos\par}
\vspace{0.12cm}
{\large\etiqueta{colorBM}{BM} \quad {\color{tonoOscuro}Andrew Blake y Haroon Mumtaz}\par}
{\normalsize \emph{\color{colorBM}Applied Bayesian Econometrics for Central Bankers}\par}
{\small\color{tonoPizarra} Bank of England, CCBS, edición actualizada, 2017 --- 7 capítulos\par}
\vspace{0.12cm}
{\large\etiqueta{colorFa}{Fa} \quad {\color{tonoOscuro}Ray C. Fair}\par}
{\normalsize \emph{\color{colorFa}Macroeconomic Modeling: The Cowles Commission Approach}\par}
{\small\color{tonoPizarra} MIT Press, 2024 --- 18 capítulos\par}
\vspace{0.12cm}
{\large\etiqueta{colorET}{ET} \quad {\color{tonoOscuro}Graham Elliott y Allan Timmermann}\par}
{\normalsize \emph{\color{colorET}Economic Forecasting}\par}
{\small\color{tonoPizarra} Princeton University Press, 2016 --- 21 capítulos\par}
\vspace{0.12cm}
{\large\etiqueta{colorPs}{Ps} \quad {\color{tonoOscuro}M. Hashem Pesaran}\par}
{\normalsize \emph{\color{colorPs}Time Series and Panel Data Econometrics}\par}
{\small\color{tonoPizarra} Oxford University Press, 2015 --- capítulo 17\par}
\vspace{0.35cm}
{\normalsize\color{colorPUCP} Con temas y lecturas de tres cursos de la PUCP (siete programas) y uno del BCRP\par}
\vfill
{\setlength{\fboxsep}{6pt}\colorbox{tonoPrincipal}{\normalsize\sffamily\bfseries\color{white} Quince bloques temáticos y un pretramo --- 10 meses}\par}
\vspace{0.2cm}
\end{titlepage}

\tableofcontents
\newpage

%------------------------------------------------------------------------------
\section*{Cómo usar este sílabo}
\addcontentsline{toc}{section}{Cómo usar este sílabo}

El plan está organizado en \textbf{quince bloques temáticos} y un pretramo, y
sigue el orden de los cursos de series de tiempo de la PUCP: primero los
procesos univariados estacionarios y no estacionarios, luego los modelos
multivariados (VAR, VAR estructurales y cointegración), después los modelos de
espacio de estados, volatilidad y no linealidad, a continuación la
econometría bayesiana y, al final, la estimación de modelos DSGE y la
construcción y evaluación de modelos macroeconométricos. Cierra con un bloque
dedicado al pronóstico económico, que reúne los modelos de todo el plan bajo
una misma pregunta: cómo pronosticar, cómo evaluar y comparar pronósticos y
cómo combinarlos. Conviene cerrar cada bloque antes de abrir el siguiente.

\medskip
\noindent\destacado{Requisito.} El plan supone el sílabo de econometría,
en particular la teoría asintótica (Bloque 4), el GMM (Bloque 9) y la máxima
verosimilitud (Bloque 10). El pretramo los repasa en el lenguaje de las series
de tiempo. Este sílabo desarrolla el módulo opcional de macroeconometría del
sílabo de macroeconomía (Anexo G.3) y lo amplía.

\medskip
\noindent\destacado{Criterio general de reparto} (en el orden en que se lee cada tema):
\begin{itemize}[leftmargin=1.4em,itemsep=2pt]
  \item \textcolor{colorEn}{\textbf{Enders}} aporta la intuición y las
        aplicaciones. Se lee \emph{primero} en los Bloques 1 a 10.
  \item \textcolor{colorMHH}{\textbf{Martin, Hurn y Harris}} es el texto base
        del tramo de series de tiempo. Trata todos los modelos con un mismo
        marco ---la máxima verosimilitud--- y fija la notación.
  \item \textcolor{colorHm}{\textbf{Hamilton}} es la referencia rigurosa:
        asintótica, raíces unitarias, cointegración, filtro de Kalman y
        cambio de régimen. Es el texto básico de 1ECO24 y de 2ECO01.
  \item \textcolor{colorKL}{\textbf{Kilian y Lütkepohl}} es la lectura
        principal de los VAR estructurales (Bloques 5 y 6) y de las
        proyecciones locales.
  \item \textcolor{colorBM}{\textbf{Blake y Mumtaz}} es la lectura principal
        de la econometría bayesiana (Bloques 11 a 13), con códigos en Matlab
        para replicar cada ejemplo.
  \item \textcolor{colorBE}{\textbf{Bårdsen et al.}} y
        \textcolor{colorFa}{\textbf{Fair}} sostienen el Bloque 14: cómo se
        construye, se evalúa y se usa un modelo macroeconométrico.
  \item \textcolor{colorET}{\textbf{Elliott y Timmermann}} es la lectura
        principal del Bloque 15 (pronóstico). En los bloques anteriores
        aparece solo como consulta; en el Bloque 15 entra en las horas.
        \textcolor{colorPs}{\textbf{Pesaran}} (cap.~17) da un resumen
        compacto de los mismos temas y sirve de profundización.
\end{itemize}

\medskip
\noindent\destacado{Qué entra en las horas.} Las líneas \emph{Lecturas},
\emph{Profundización}, \emph{Artículos} y \emph{Artículos (PUCP)} entran en el
cómputo de horas.
Las líneas \emph{Consulta} y \emph{Consulta (PUCP)} remiten a libros de la
carpeta y a referencias de los programas que no entran (Anexo E).

\medskip
\noindent\destacado{Material excluido.} Quedan fuera las partes de Martin, Hurn y
Harris que repiten el sílabo de econometría (regresión lineal y no lineal,
métodos no paramétricos, datos discretos), los temas propios de finanzas
(valor en riesgo, CoVaR) y el panel VAR, que pasan al material opcional. El
detalle está en el Anexo B.

\medskip
\noindent\destacado{Sobre el software y los ejercicios.} La macroeconometría se
aprende estimando. Enders trae sus bases de datos en la carpeta (en formatos
EViews, SAS y texto); Blake y Mumtaz, códigos en Matlab para cada capítulo;
Martin, Hurn y Harris, ejercicios con programas en el sitio del libro.
Conviene replicar cada semana al menos una aplicación, de preferencia con
datos del BCRP.

\medskip
\noindent\destacado{Notación.} \MHH{n} remite al capítulo $n$ de Martin, Hurn y
Harris, y \BE{n}, \EN{n}, \HM{n}, \KL{n}, \BM{n}, \FA{n}, \ET{n} y \PS{n}, al
capítulo de los demás textos. Si solo entra parte de un capítulo se indican secciones:
\HM{17.6} es la sección 6 del capítulo 17 de Hamilton. Enders y Blake y Mumtaz
numeran las secciones dentro de cada capítulo: \EN{5.13} es la sección 13 del
capítulo 5. \emph{Consulta} remite a libros de la carpeta que no entran en el
cómputo de horas (Anexo E).

%------------------------------------------------------------------------------
\newpage
\section*{Cronograma por semanas}
\addcontentsline{toc}{section}{Cronograma por semanas}

Calendario del plan: 430 horas distribuidas en \textbf{43 semanas de
contenido}, que equivalen a \textbf{unos 10 meses de calendario}, a un ritmo de
diez horas semanales.

\medskip
\noindent Todas las semanas son de materia nueva y no hay semanas de colchón.
La columna de horas es la que gobierna el reparto. Si hay semanas perdidas por
feriados, viajes o imprevistos, el ajuste correcto es alargar el calendario,
no comprimir los bloques.

\medskip
{\filasalternas
\noindent\begin{tabularx}{\textwidth}{lXcc}
\toprule
\rowcolor{tonoOscuro}\cab{Semanas} & \cab{Contenido} & \cab{Sem.} & \cab{Horas} \\
\midrule
1 & Pretramo --- Máxima verosimilitud para series de tiempo & 1 & 10 \\
2--3 & Bloque 1 --- Ecuaciones en diferencias y estacionariedad & 2 & 20 \\
\rowcolor{tonoHielo}\textbf{4--6} & \textbf{Bloque 2 --- Modelos ARMA y pronóstico} & \textbf{3} & \textbf{30} \\
7--8 & Bloque 3 --- Asintótica, regresión dinámica y análisis espectral & 2 & 20 \\
\rowcolor{tonoHielo}\textbf{9--11} & \textbf{Bloque 4 --- Tendencias y raíces unitarias} & \textbf{3} & \textbf{30} \\
\rowcolor{tonoHielo}\textbf{12--14} & \textbf{Bloque 5 --- Vectores autorregresivos} & \textbf{3} & \textbf{30} \\
\rowcolor{tonoHielo}\textbf{15--18} & \textbf{Bloque 6 --- VAR estructurales y proyecciones locales} & \textbf{4} & \textbf{40} \\
\rowcolor{tonoHielo}\textbf{19--21} & \textbf{Bloque 7 --- Cointegración y modelos VECM} & \textbf{3} & \textbf{30} \\
\rowcolor{tonoHielo}\textbf{22--24} & \textbf{Bloque 8 --- Espacio de estados, filtro de Kalman y factores} & \textbf{3} & \textbf{30} \\
25--26 & Bloque 9 --- Modelos de volatilidad & 2 & 20 \\
27--28 & Bloque 10 --- Modelos no lineales y regresión cuantílica & 2 & 20 \\
\rowcolor{tonoHielo}\textbf{29--31} & \textbf{Bloque 11 --- Econometría bayesiana: fundamentos y MCMC} & \textbf{3} & \textbf{30} \\
\rowcolor{tonoHielo}\textbf{32--34} & \textbf{Bloque 12 --- VAR bayesianos y modelos con parámetros cambiantes} & \textbf{3} & \textbf{30} \\
35--36 & Bloque 13 --- Estimación de modelos DSGE & 2 & 20 \\
\rowcolor{tonoHielo}\textbf{37--39} & \textbf{Bloque 14 --- Modelización macroeconométrica y pronóstico} & \textbf{3} & \textbf{30} \\
\rowcolor{tonoHielo}\textbf{40--43} & \textbf{Bloque 15 --- Pronóstico económico: evaluación, combinación y \emph{nowcasting}} & \textbf{4} & \textbf{40} \\
\midrule
\rowcolor{tonoHielo}\multicolumn{2}{l}{\textbf{Total}} & \textbf{43} & \textbf{430} \\
\bottomrule
\end{tabularx}
}

\medskip
\noindent\destacado{Ritmo supuesto.} Diez horas por semana, todas de contenido
nuevo: por ejemplo, cuatro sesiones de dos horas entre semana para la teoría y
una sesión de dos horas el sábado para estimar con datos.

\medskip
\noindent\destacado{Los tramos largos.} Los Bloques 5 a 7 (semanas 12--21) forman
el núcleo del plan: el análisis multivariado que usan los bancos centrales.
Los Bloques 11 a 13 (semanas 29--36) forman un recorrido continuo por la
econometría bayesiana y son el punto donde el avance suele atascarse, porque
exigen programar. El Bloque 15 (semanas 40--43) vuelve sobre todos los
modelos anteriores desde el punto de vista del pronóstico y se apoya, sobre
todo, en los Bloques 2, 5, 8 y 12.

\medskip
\noindent\destacado{Hitos de control.} Conviene verificar el avance al terminar las
semanas 11, 21, 28, 36 y 39. Si en alguno de esos cortes se acumula más de una
semana de retraso, se alarga el calendario y no se comprime el bloque
siguiente.

\newpage
\subsection*{Detalle semana a semana}

{\footnotesize\filasalternas
\renewcommand{\arraystretch}{1.1}\noindent\begin{tabularx}{\textwidth}{clXc}
\toprule
\rowcolor{tonoOscuro}\cab{Sem.} & \cab{Bloque} & \cab{Temas} & \cab{Horas} \\
\midrule
1  & Pretramo & P.1 La verosimilitud en series de tiempo; P.2 Optimización numérica y pruebas clásicas & 10 \\
\midrule
2  & 1 & 1.1 Las series macroeconómicas; 1.2 Ecuaciones en diferencias & 10 \\
3  & 1 & 1.3 Operadores de rezagos; 1.4 Estacionariedad y ergodicidad & 10 \\
\midrule
4  & 2 & 2.1 Procesos MA, AR y ARMA & 10 \\
5  & 2 & 2.2 Metodología de Box--Jenkins; 2.3 Estimación por MCO y máxima verosimilitud & 10 \\
6  & 2 & 2.4 Pronóstico; 2.5 Estacionalidad e inestabilidad de parámetros & 10 \\
\midrule
7  & 3 & 3.1 Teoremas límite con dependencia; 3.2 Regresión dinámica y modelos ADL & 10 \\
8  & 3 & 3.3 Análisis espectral; 3.4 Filtros lineales: Hodrick--Prescott y paso de banda & 10 \\
\midrule
9  & 4 & 4.1 Tendencias determinísticas y estocásticas; 4.2 Movimiento browniano y TLC funcional & 10 \\
10 & 4 & 4.3 Pruebas de Dickey--Fuller y Phillips--Perron; 4.4 Potencia y pruebas más potentes & 10 \\
11 & 4 & 4.5 Quiebres estructurales; 4.6 Descomposiciones de tendencia y ciclo & 10 \\
\midrule
12 & 5 & 5.1 Representación de un VAR; 5.2 Estimación y especificación & 10 \\
13 & 5 & 5.3 Causalidad de Granger y pronóstico & 10 \\
14 & 5 & 5.4 Funciones de impulso-respuesta y descomposición de varianza & 10 \\
\midrule
15 & 6 & 6.1 De la forma reducida a la estructural; 6.2 Restricciones de corto plazo & 10 \\
16 & 6 & 6.3 Restricciones de largo plazo; 6.4 Inferencia y bootstrap & 10 \\
17 & 6 & 6.5 Restricciones de signos; 6.6 Instrumentos externos y heteroscedasticidad & 10 \\
18 & 6 & 6.7 Proyecciones locales; 6.8 Descomposiciones históricas y contrafactuales & 10 \\
\midrule
19 & 7 & 7.1 Regresión espuria; 7.2 Cointegración y método de Engle--Granger & 10 \\
20 & 7 & 7.3 El VECM y el método de Johansen & 10 \\
21 & 7 & 7.4 Identificación, pruebas y dinámica; 7.5 VAR estructurales con cointegración & 10 \\
\midrule
22 & 8 & 8.1 Forma de espacio de estados; 8.2 Filtro de Kalman y suavizado & 10 \\
23 & 8 & 8.3 Estimación por máxima verosimilitud; 8.4 Parámetros cambiantes y variables no observables & 10 \\
24 & 8 & 8.5 Modelos de factores, componentes principales y FAVAR & 10 \\
\midrule
25 & 9 & 9.1 Modelos ARCH y GARCH; 9.2 Extensiones univariadas & 10 \\
26 & 9 & 9.3 GARCH multivariado & 10 \\
\midrule
27 & 10 & 10.1 Modelos de umbral y de transición suave; 10.2 Cambio de régimen de Markov & 10 \\
28 & 10 & 10.3 Dinámica no lineal y redes neuronales; 10.4 Regresión cuantílica y crecimiento en riesgo & 10 \\
\midrule
29 & 11 & 11.1 El enfoque bayesiano; 11.2 Regresión lineal bayesiana & 10 \\
30 & 11 & 11.3 Muestreo de Gibbs & 10 \\
31 & 11 & 11.4 Metropolis--Hastings, convergencia y Monte Carlo hamiltoniano & 10 \\
\midrule
32 & 12 & 12.1 VAR bayesianos y sus distribuciones a priori & 10 \\
33 & 12 & 12.2 VAR bayesianos estructurales y pronóstico condicional; 12.3 Espacio de estados bayesiano & 10 \\
34 & 12 & 12.4 TVP-VAR y volatilidad estocástica; 12.5 FAVAR, frecuencias mixtas y cambio de régimen & 10 \\
\midrule
35 & 13 & 13.1 VAR y DSGE; 13.2 Estimación bayesiana de modelos DSGE & 10 \\
36 & 13 & 13.3 GMM, calibración y métodos de simulación & 10 \\
\midrule
37 & 14 & 14.1 Metodología de la modelización; 14.2 Exogeneidad, invarianza y crítica de Lucas & 10 \\
38 & 14 & 14.3 Inflación, salarios y curva de Phillips; 14.4 La curva de Phillips neokeynesiana y el dinero & 10 \\
39 & 14 & 14.5 Transmisión monetaria, reglas y multiplicadores; 14.6 Pronóstico y su evaluación & 10 \\
\midrule
40 & 15 & 15.1 Pérdida y pronóstico óptimo; 15.2 Diseño fuera de muestra y datos en tiempo real & 10 \\
41 & 15 & 15.3 Selección de modelos y pronóstico univariado; 15.4 Pronóstico multivariado & 10 \\
42 & 15 & 15.5 Evaluación y comparación de pronósticos; 15.6 Combinación e inestabilidad & 10 \\
43 & 15 & 15.7 Densidades, intervalos y \emph{fan charts}; 15.8 \emph{Nowcasting} y frecuencias mixtas & 10 \\
\bottomrule
\end{tabularx}
}

%------------------------------------------------------------------------------
\newpage
\bloque{Pretramo --- Máxima verosimilitud para series de tiempo --- 1 semana}
{Repaso del principio de máxima verosimilitud en el marco que usa Martin, Hurn y
Harris para todo el libro. Supone el Bloque 10 del sílabo de econometría; aquí
se vuelve a ver con observaciones dependientes y con la notación del texto
base.}

\tema{P.1 La verosimilitud en series de tiempo}
\lecturas{\MHH{1} \quad \MHH{2.1--2.5}}
\nota{Distribuciones conjuntas y condicionales de observaciones dependientes,
descomposición de la verosimilitud por predicción, condiciones de regularidad
y propiedades asintóticas.}

\tema{P.2 Optimización numérica y pruebas clásicas}
\lecturas{\MHH{3} \quad \MHH{4}}
\apoyo{\HM{5.7}}
\nota{Métodos de Newton y cuasi-Newton, cálculo de errores estándar y pruebas
de razón de verosimilitud, Wald y multiplicador de Lagrange.}

%------------------------------------------------------------------------------
\bloque{Bloque 1 --- Ecuaciones en diferencias y estacionariedad --- 2 semanas}
{Los temas 1 a 3 de 1ECO24 (2026-1) y la sesión 1.1 de 1ECO78: las
herramientas para describir la dinámica de una serie antes de estimarla.}

\tema{1.1 Las series de tiempo macroeconómicas}
\lecturas{\EN{1.1} \quad \MHH{13.1--13.2}}
\consultalib{Fair (2024), cap.~2 \quad Neusser (2016), cap.~1}
\nota{Hechos estilizados de las series macroeconómicas y financieras:
tendencia, persistencia, estacionalidad y volatilidad cambiante.}

\tema{1.2 Ecuaciones en diferencias y sus soluciones}
\lecturas{\EN{1.2--1.8}}
\apoyo{\HM{1}}
\nota{Solución por iteración, ecuaciones homogéneas, soluciones particulares y
método de coeficientes indeterminados. Multiplicadores dinámicos.}

\tema{1.3 Operadores de rezagos}
\lecturas{\EN{1.9}}
\apoyo{\HM{2}}

\tema{1.4 Estacionariedad, ergodicidad y ruido blanco}
\lecturas{\EN{2.1} \quad \EN{2.3} \quad \MHH{13.4}}
\apoyo{\HM{3.1--3.2}}
\consultalib{Brockwell y Davis (2009), cap.~1}

%------------------------------------------------------------------------------
\bloque{Bloque 2 --- Modelos ARMA y pronóstico --- 3 semanas}
{El análisis univariado estacionario: el tema 4 de 1ECO24, la sesión 1.1 de
1ECO78 y el tema 1 de la Unidad 4 de 2ECO01.}

\tema{2.1 Procesos MA, AR y ARMA}
\lecturas{\EN{2.2} \quad \EN{2.4--2.7} \quad \MHH{13.3} \quad \MHH{13.5}}
\apoyo{\HM{3.3--3.7}}
\consultalib{Neusser (2016), cap.~2 \quad Shumway y Stoffer, cap.~3}
\nota{Condiciones de estacionariedad e invertibilidad, funciones de
autocorrelación simple y parcial, función generadora de autocovarianzas.}

\tema{2.2 Metodología de Box--Jenkins}
\lecturas{\EN{2.8} \quad \EN{2.10} \quad \MHH{13.7} \quad \MHH{13.9}}
\apoyo{\HM{4.8}}
\nota{Identificación, elección del número de rezagos con criterios de
información, estimación y diagnóstico de residuos.}

\tema{2.3 Estimación por MCO y por máxima verosimilitud}
\lecturas{\MHH{13.6} \quad \MHH{13.8}}
\apoyo{\HM{5.1--5.6} \quad \HM{5.8}}
\nota{Verosimilitud exacta y condicional de procesos AR, MA y ARMA gaussianos.}

\tema{2.4 Pronóstico}
\lecturas{\EN{2.9} \quad \EN{2.13}}
\apoyo{\HM{4.1--4.3} \quad \HM{4.6}}
\consultalib{Elliott y Timmermann (2016), caps.~2 y 7}
\nota{Pronóstico óptimo bajo pérdida cuadrática, error de predicción e
intervalos de predicción, y descomposición de Wold. La evaluación, la
comparación y la combinación de pronósticos se estudian en el Bloque 15.}

\tema{2.5 Estacionalidad e inestabilidad de parámetros}
\lecturas{\EN{2.11--2.12}}

%------------------------------------------------------------------------------
\bloque{Bloque 3 --- Asintótica, regresión dinámica y análisis espectral --- 2 semanas}
{La teoría asintótica con observaciones dependientes (tema 7 de 1ECO24 y
sesión 2.1 de 1ECO78), la regresión con series de tiempo y el análisis en el
dominio de la frecuencia (sesión 2.4 de 1ECO78).}

\tema{3.1 Teoremas límite para observaciones dependientes}
\lecturas{\HM{7}}
\consultalib{White (2001), caps.~3 y 5 \quad Hayashi (2000), cap.~6 \quad Hansen (2021), cap.~14}
\nota{Ergodicidad, diferencias martingala, ley de los grandes números y teorema
central del límite para procesos dependientes; estimación de la varianza de
largo plazo.}

\tema{3.2 Regresión dinámica y modelos autorregresivos de rezagos distribuidos}
\lecturas{\EN{5.1--5.4} \quad \MHH{7.1--7.7}}
\apoyo{\HM{8}}
\consultalib{Davidson y MacKinnon (2004), cap.~13.4}
\nota{Funciones de transferencia, análisis de intervención y regresiones con
errores autocorrelacionados y variables dependientes rezagadas.}

\tema{3.3 Análisis espectral}
\lecturas{\HM{6.1--6.3}}
\apoyo{\HM{10.4}}
\consultalib{Neusser (2016), cap.~6 \quad Brockwell y Davis (2009), caps.~4 y 10}
\nota{El espectro poblacional, el periodograma y su estimación.}

\tema{3.4 Filtros lineales: Hodrick--Prescott y paso de banda}
\lecturas{\HM{6.4}}
\consulta{Christiano y Fitzgerald (2003) \quad Cogley y Nason (1995)}
\nota{Qué hace cada filtro en el dominio de la frecuencia y por qué el filtro
de Hodrick--Prescott puede generar ciclos espurios. Prepara la discusión del
producto potencial en el tema 8.4.}

%------------------------------------------------------------------------------
\bloque{Bloque 4 --- Tendencias y raíces unitarias --- 3 semanas}
{La no estacionariedad univariada: temas 8 y 9 de 1ECO24, sesiones 2.1 y 2.2
de 1ECO78 y tema 3 de la Unidad 4 de 2ECO01. La asintótica no estándar del
tema 4.2 es la parte más técnica del plan.}

\tema{4.1 Tendencias determinísticas y estocásticas}
\lecturas{\EN{4.1--4.3} \quad \MHH{16.1--16.3}}
\apoyo{\HM{15} \quad \HM{16}}

\tema{4.2 Movimiento browniano, teorema central del límite funcional y asintótica no estándar}
\lecturas{\MHH{16.4--16.5}}
\apoyo{\HM{17.1--17.5} \quad \HM{18.1}}
\consultalib{White (2001), cap.~7}

\tema{4.3 Pruebas de Dickey--Fuller, Dickey--Fuller aumentada y Phillips--Perron}
\lecturas{\EN{4.4--4.7} \quad \MHH{17.1--17.5} \quad \MHH{17.7}}
\apoyo{\HM{17.6--17.7}}
\consultalib{Pfaff (2008) \quad Hayashi (2000), cap.~9}

\tema{4.4 Potencia, regresores determinísticos y pruebas más potentes}
\lecturas{\EN{4.9--4.10} \quad \MHH{17.6}}
\apoyo{\HM{17.8}}
\consulta{Elliott, Rothenberg y Stock (1996)}
\nota{Pruebas DF-GLS, KPSS y otras alternativas.}

\tema{4.5 Quiebres estructurales y raíces unitarias}
\lecturas{\EN{4.8} \quad \EN{7.12} \quad \MHH{17.8}}
\consultalib{Kilian y Lütkepohl (2017), cap.~19}

\tema{4.6 Descomposiciones de tendencia y ciclo}
\lecturas{\EN{4.12}}
\consultalib{Neusser (2016), apéndice D}
\nota{Descomposición de Beveridge--Nelson y modelos de componentes no
observados; se completan con el filtro de Kalman en el tema 8.4.}

%------------------------------------------------------------------------------
\bloque{Bloque 5 --- Vectores autorregresivos --- 3 semanas}
{El tema 10 de 1ECO24, la sesión 1.2 de 1ECO78 y el tema 2 de la Unidad 4 de
2ECO01. Kilian y Lütkepohl es la lectura principal desde aquí hasta el
Bloque 7.}

\tema{5.1 Representación de un VAR: estabilidad y representación de Wold}
\lecturas{\EN{5.5} \quad \KL{2.1--2.2} \quad \MHH{13.10}}
\apoyo{\HM{10.1--10.3}}
\consultalib{Lütkepohl (2005), cap.~2 \quad Neusser (2016), caps.~10 y 12}

\tema{5.2 Estimación, selección de rezagos y diagnóstico}
\lecturas{\EN{5.6} \quad \KL{2.3} \quad \KL{2.6--2.8}}
\apoyo{\HM{11.1} \quad \HM{11.3}}
\consultalib{Lütkepohl (2005), caps.~3--5}

\tema{5.3 Causalidad de Granger y pronóstico}
\lecturas{\EN{5.8} \quad \KL{2.4--2.5} \quad \KL{7.1--7.4}}
\apoyo{\HM{11.2}}
\consultalib{Elliott y Timmermann (2016), cap.~9}
\nota{Kilian y Lütkepohl explican por qué la causalidad de Granger dejó de
usarse como prueba de causalidad en macroeconomía.}

\tema{5.4 Funciones de impulso-respuesta y descomposición de la varianza}
\lecturas{\EN{5.7} \quad \KL{4.1--4.2}}
\apoyo{\HM{11.4--11.5}}
\articulos{Sims (1980) \quad Stock y Watson (2001)}

%------------------------------------------------------------------------------
\bloque{Bloque 6 --- VAR estructurales y proyecciones locales --- 4 semanas}
{El tema 11 de 1ECO24, la sesión 1.3 de 1ECO78 y el VAR estructural del curso
del BCRP. Es el bloque más largo: la identificación de choques es la
herramienta central del análisis de política monetaria y fiscal.}

\tema{6.1 De la forma reducida a la estructural: choques estructurales}
\lecturas{\EN{5.10} \quad \KL{7.5--7.6} \quad \KL{8.1--8.3} \quad \MHH{14.1--14.3}}
\apoyo{\HM{11.6}}

\tema{6.2 Identificación por restricciones de corto plazo}
\lecturas{\EN{5.11--5.12} \quad \KL{8.4--8.5} \quad \KL{9.1--9.4} \quad \MHH{14.4--14.5}}
\articulos{Bernanke y Mihov (1998)}
\nota{Modelos recursivos (Cholesky) y no recursivos; estimación por método de
momentos, variables instrumentales y máxima verosimilitud con información
completa.}

\tema{6.3 Identificación por restricciones de largo plazo}
\lecturas{\EN{5.13--5.14} \quad \KL{10} \quad \KL{11.1--11.3}}
\articulos{Blanchard y Quah (1989)}

\tema{6.4 Inferencia: método delta y bootstrap}
\lecturas{\KL{12.1--12.5}}
\apoyo{\HM{11.7}}
\nota{Bandas de confianza para funciones de impulso-respuesta, bootstrap con
corrección de sesgo y errores frecuentes en la inferencia.}

\tema{6.5 Identificación por restricciones de signos}
\lecturas{\KL{13.1--13.6} \quad \KL{13.9}}
\articulos{Uhlig (2005)}

\tema{6.6 Instrumentos externos, datos de alta frecuencia y heteroscedasticidad}
\lecturas{\KL{14} \quad \KL{15}}
\nota{Identificación con instrumentos externos (\emph{proxy} SVAR), con
sorpresas de política medidas en futuros de alta frecuencia y con cambios de
volatilidad.}

\tema{6.7 Proyecciones locales}
\lecturas{\KL{12.8} \quad Montiel Olea, Plagborg-Møller, Qian y Wolf (2025)}
\articulos{Jordà (2005)}
\nota{Proyecciones locales frente a VAR: mismo objeto de estimación, distinto
equilibrio entre sesgo y varianza. El artículo de Montiel Olea et al. está en
la carpeta.}

\tema{6.8 Descomposiciones históricas, escenarios y contrafactuales}
\lecturas{\KL{4.3--4.6}}

%------------------------------------------------------------------------------
\bloque{Bloque 7 --- Cointegración y modelos VECM --- 3 semanas}
{La no estacionariedad multivariada: el tema 12 de 1ECO24, la sesión 2.3 de
1ECO78 y el tema 3 de la Unidad 4 de 2ECO01. Martin, Hurn y Harris es la
lectura principal.}

\tema{7.1 Variables integradas y regresión espuria}
\lecturas{\EN{6.1--6.2}}
\apoyo{\HM{18.2--18.3}}

\tema{7.2 Cointegración, corrección de errores y método de Engle--Granger}
\lecturas{\EN{6.3--6.6} \quad \MHH{18.1--18.2}}
\apoyo{\HM{19.1--19.3}}
\consultalib{Pfaff (2008)}

\tema{7.3 El VECM y el método de máxima verosimilitud de Johansen}
\lecturas{\EN{6.7--6.9} \quad \MHH{18.3--18.4} \quad \KL{3.1--3.3}}
\apoyo{\HM{20.1--20.3}}
\consultalib{Lütkepohl (2005), caps.~6--8}

\tema{7.4 Identificación, pruebas y dinámica de los sistemas cointegrados}
\lecturas{\EN{6.10--6.11} \quad \MHH{18.5--18.8} \quad \KL{3.4--3.6}}
\apoyo{\HM{20.4}}
\nota{Identificación de los vectores de cointegración, pruebas sobre ellos y
funciones de impulso-respuesta en el VECM.}

\tema{7.5 VAR estructurales con variables cointegradas}
\lecturas{\KL{12.6--12.7}}
\consultalib{Lütkepohl (2005), cap.~9 \quad Garratt, Lee, Pesaran y Shin (2006)}
\nota{Inferencia en VAR integrados y cerca de la raíz unitaria; el enfoque
estructural de largo plazo de Garratt et al. queda como consulta.}

%------------------------------------------------------------------------------
\bloque{Bloque 8 --- Espacio de estados, filtro de Kalman y modelos de factores --- 3 semanas}
{El tema 13 de 1ECO24, las sesiones 1.4 y 2.4 de 1ECO78 y el curso del BCRP.
La forma de espacio de estados es el puente hacia los modelos bayesianos con
parámetros cambiantes (Bloque 12) y hacia la estimación de modelos DSGE
(Bloque 13).}

\tema{8.1 La forma de espacio de estados}
\lecturas{\MHH{15.1--15.2}}
\apoyo{\HM{13.1}}
\consultalib{Lütkepohl (2005), cap.~18 \quad Neusser (2016), cap.~17}

\tema{8.2 El filtro de Kalman: predicción, actualización y suavizado}
\lecturas{\MHH{15.3--15.4} \quad \BM{3.4}}
\apoyo{\HM{13.2--13.3} \quad \HM{13.5--13.6}}

\tema{8.3 Estimación por máxima verosimilitud e inferencia}
\lecturas{\MHH{15.6}}
\apoyo{\HM{13.4} \quad \HM{13.7}}

\tema{8.4 Parámetros cambiantes y variables no observables}
\lecturas{\HM{13.8}}
\consulta{Kim y Nelson (1999) \quad Harvey (1990)}
\nota{Producto potencial, tasa natural y tendencias como estados no
observables; el filtro de Hodrick--Prescott como caso particular del
suavizado de Kalman.}

\tema{8.5 Modelos de factores, componentes principales y FAVAR}
\lecturas{\MHH{15.5} \quad \MHH{15.7} \quad \KL{16.1--16.2}}
\consultalib{Hansen (2021), caps.~11.12--11.16}
\articulos{Stock y Watson (2002)}
\nota{Extracción de factores con el filtro de Kalman y con componentes
principales; relación con los modelos VARMA. Es el tema 9 de 1ECO24
(2025-2).}

%------------------------------------------------------------------------------
\bloque{Bloque 9 --- Modelos de volatilidad --- 2 semanas}
{El tema 14 de 1ECO24 y el módulo de volatilidad de 1ECO78. La volatilidad
estocástica se estudia en el Bloque 12, con métodos bayesianos.}

\tema{9.1 Modelos ARCH y GARCH}
\lecturas{\EN{3.1--3.5} \quad \EN{3.8} \quad \MHH{20.1--20.3}}
\apoyo{\HM{21.1} \quad \MHH{9.1--9.4}}
\consultalib{Tsay (2010), cap.~3}
\nota{Estimación por máxima verosimilitud y por cuasi-máxima verosimilitud
(\MHH{9}).}

\tema{9.2 Extensiones univariadas}
\lecturas{\EN{3.6--3.7} \quad \EN{3.9--3.10} \quad \MHH{20.4--20.5}}
\apoyo{\HM{21.2}}
\nota{GARCH-M, EGARCH, TARCH, IGARCH y FIGARCH; distribuciones condicionales
no normales.}

\tema{9.3 GARCH multivariado e impulso-respuesta de la volatilidad}
\lecturas{\EN{3.11--3.12} \quad \MHH{20.6}}
\consultalib{Lütkepohl (2005), cap.~16 \quad Tsay (2013), cap.~7}
\consulta{Engle (2002)}

%------------------------------------------------------------------------------
\bloque{Bloque 10 --- Modelos no lineales y regresión cuantílica --- 2 semanas}
{Los temas opcionales de 1ECO24 (umbrales y cambio de régimen), los modelos no
lineales de 1ECO78, el curso del BCRP y la regresión cuantílica con series de
tiempo (sesión 2.4 de 1ECO78).}

\tema{10.1 Pruebas de no linealidad, modelos de umbral y de transición suave}
\lecturas{\EN{7.1--7.7} \quad \EN{7.9} \quad \MHH{19.1--19.3}}
\consultalib{Chen y Tsay (2019)}
\nota{Modelos TAR, SETAR, LSTAR y ESTAR (las regresiones de umbral y de
transición suave del curso del BCRP).}

\tema{10.2 Cambio de régimen de Markov}
\lecturas{\EN{7.8} \quad \MHH{19.6}}
\apoyo{\HM{22}}
\consulta{Kim y Nelson (1999)}
\nota{El filtro de Hamilton; la versión bayesiana se ve en el tema 12.5.}

\tema{10.3 Dinámica no lineal, redes neuronales y VAR no lineales}
\lecturas{\EN{7.10--7.11} \quad \MHH{19.4} \quad \MHH{19.8} \quad \KL{18.1--18.4}}
\nota{Funciones de impulso-respuesta generalizadas, raíces unitarias con
ajuste no lineal y redes neuronales artificiales, la introducción al
aprendizaje automático que piden 1ECO24 (2025-2) y 1ECO78.}

\tema{10.4 Regresión cuantílica con series de tiempo y crecimiento en riesgo}
\lecturas{Hansen (2021), cap.~24}
\articulos{Adrian, Boyarchenko y Giannone (2019)}
\consultalib{Elliott y Timmermann (2016), cap.~13}
\nota{Pronóstico de la densidad condicional del PBI y medición del
\emph{growth at risk}. Hansen está en la carpeta de econometría y es el texto
del tema 13.5 de ese sílabo. La evaluación de estas densidades se ve en el
tema 15.7.}

%------------------------------------------------------------------------------
\bloque{Bloque 11 --- Econometría bayesiana: fundamentos y MCMC --- 3 semanas}
{\textbf{Bloque nuevo respecto de 1ECO24.} Es la sesión 3.1 y el módulo 4 de
1ECO78. Blake y Mumtaz es la lectura principal y trae los códigos para
replicar cada algoritmo; Hamilton y Greene dan la base formal.}

\tema{11.1 El enfoque bayesiano: distribución a priori, verosimilitud y posterior}
\lecturas{\BM{1.1--1.2}}
\apoyo{\HM{12.1}}
\consultalib{Greene (2012), caps.~16.1--16.4 \quad McElreath (2016), caps.~1--3}
\nota{El teorema de Bayes, las distribuciones marginales y condicionales y las
diferencias entre la econometría clásica y la bayesiana.}

\tema{11.2 El modelo de regresión lineal bayesiano}
\lecturas{\BM{1.2} \quad \BM{1.5}}
\consultalib{Greene (2012), cap.~16.3 \quad McElreath (2016), caps.~4--6}
\consulta{Koop (2003)}
\nota{Distribuciones a priori conjugadas e informativas, elección de
hiperparámetros, verosimilitud marginal y comparación de modelos.}

\tema{11.3 El muestreo de Gibbs}
\lecturas{\BM{1.3}}
\apoyo{\HM{12.3}}
\consultalib{Greene (2012), cap.~16.5}
\articulos{Casella y George (1992)}
\nota{Errores autocorrelacionados y heteroscedásticos como extensiones del
muestreador.}

\tema{11.4 Metropolis--Hastings, convergencia y Monte Carlo hamiltoniano}
\lecturas{\BM{5.1--5.4} \quad \BM{5.6} \quad \BM{5.8}}
\consultalib{McElreath (2016), cap.~8}
\consulta{Chib (2001) \quad Gelman et al. (2013)}
\nota{Algoritmos de paseo aleatorio y de independencia, diagnósticos de
convergencia y cálculo de la verosimilitud marginal. El Monte Carlo
hamiltoniano y los métodos variacionales que menciona 1ECO78 se introducen con
McElreath.}

%------------------------------------------------------------------------------
\bloque{Bloque 12 --- VAR bayesianos y modelos con parámetros cambiantes --- 3 semanas}
{Las sesiones 3.1 a 3.3 de 1ECO78: los modelos que hoy usan los bancos
centrales para pronosticar y analizar la política monetaria.}

\tema{12.1 VAR bayesianos y sus distribuciones a priori}
\lecturas{\BM{2.1--2.5} \quad \BM{2.9} \quad \KL{5}}
\apoyo{\HM{12.2}}
\consultalib{Elliott y Timmermann (2016), cap.~5}
\articulos{Litterman (1986) \quad Koop y Korobilis (2010)}
\nota{Distribuciones a priori de Minnesota, normal-Wishart inversa y de estado
estacionario; observaciones ficticias. 1ECO78 añade el tratamiento de los
valores extremos de la pandemia de 2020.}

\tema{12.2 VAR bayesianos estructurales y pronóstico condicional}
\lecturas{\BM{2.6--2.7} \quad \KL{9.5} \quad \KL{13.6--13.7}}
\nota{Restricciones de signos con muestreo de Gibbs, el papel de la
distribución a priori de la matriz de rotación, pronósticos condicionales y
escenarios.}

\tema{12.3 Modelos de espacio de estados bayesianos}
\lecturas{\BM{3.1--3.5}}
\consulta{Carter y Kohn (1994) \quad Durbin y Koopman (2002)}
\nota{El algoritmo de Carter y Kohn para muestrear los estados.}

\tema{12.4 VAR con parámetros cambiantes y volatilidad estocástica}
\lecturas{\BM{3.6} \quad \BM{5.5} \quad \BM{7} \quad \KL{18.5}}
\consulta{Cogley y Sargent (2005) \quad Kim, Shephard y Chib (1998)}
\nota{El TVP-VAR-SV, la volatilidad estocástica y el modelo de factores con
parámetros cambiantes.}

\tema{12.5 FAVAR bayesiano, frecuencias mixtas, VAR grandes y cambio de régimen}
\lecturas{\BM{3.7--3.8} \quad \BM{4} \quad \KL{16.3}}
\consultalib{Elliott y Timmermann (2016), cap.~10}
\nota{\emph{Nowcasting} con datos de frecuencia mixta y modelos de cambio de
régimen de Markov estimados con el muestreo de Gibbs. Las versiones clásicas
del \emph{nowcasting} (MIDAS, ecuaciones puente y factores dinámicos con
datos faltantes) se ven en el tema 15.8.}

%------------------------------------------------------------------------------
\bloque{Bloque 13 --- Estimación de modelos DSGE --- 2 semanas}
{La sesión 3.4 de 1ECO78 y los temas de GMM y calibración del curso del BCRP.
Enlaza con el Bloque 12 del sílabo de macroeconomía.}

\tema{13.1 VAR y modelos DSGE: solución y forma de espacio de estados}
\lecturas{\KL{6} \quad \BE{Apéndice A.2}}
\consultalib{Fair (2024), cap.~16 \quad Favero (2001), cap.~8}
\nota{Solución de modelos con expectativas racionales, su forma de espacio de
estados y la relación entre los VAR y los DSGE.}

\tema{13.2 Estimación bayesiana de modelos DSGE lineales}
\lecturas{\BM{6}}
\consulta{Herbst y Schorfheide (2016) \quad Fernández-Villaverde y Rubio-Ramírez (2007)}
\nota{Elección de observables y de distribuciones a priori, el algoritmo de
Metropolis--Hastings y el diagnóstico de la estimación. Los filtros no
lineales (Kalman extendido y de partículas) quedan como consulta.}

\tema{13.3 GMM, calibración y métodos de simulación}
\lecturas{\MHH{10} \quad \MHH{12.1--12.5}}
\apoyo{\HM{14}}
\consultalib{Favero (2001), caps.~7--8}
\nota{GMM para ecuaciones de Euler, inferencia indirecta, método eficiente de
momentos y momentos simulados; la calibración como alternativa a la
estimación.}

%------------------------------------------------------------------------------
\bloque{Bloque 14 --- Modelización macroeconométrica y pronóstico --- 3 semanas}
{Los textos de modelización: cómo se construye un modelo macroeconométrico
completo, cómo se evalúa y cómo se usa para analizar la política y
pronosticar. Retoma la cointegración (Bloque 7) y la crítica de Lucas, y
prepara el Bloque 15.}

\tema{14.1 Metodología: Cowles Commission, LSE y VAR}
\lecturas{\BE{1--2} \quad \FA{1} \quad \FA{3--4} \quad \KL{6.4}}
\consultalib{Fair (2024), caps.~5--8 \quad Favero (2001), caps.~4--6}
\nota{Modelos grandes y pequeños, el papel de la teoría y de los datos, la
identificación de subsistemas. El modelo de Fair para Estados Unidos ilustra
el enfoque de la Cowles Commission.}

\tema{14.2 Exogeneidad, invarianza y crítica de Lucas}
\lecturas{\BE{4.5} \quad \BE{9.4} \quad \BE{Apéndice A.1}}

\tema{14.3 Inflación, salarios y curva de Phillips}
\lecturas{\BE{3--6}}
\nota{El modelo \emph{main-course} de inflación en economías abiertas, la
curva de Phillips y la NAIRU, la negociación salarial y la dinámica de
salarios y precios con corrección de errores.}

\tema{14.4 La curva de Phillips neokeynesiana y el dinero}
\lecturas{\BE{7--8}}
\nota{Estimación y evaluación de la NKPC; demanda de dinero e inflación.}

\tema{14.5 Transmisión monetaria, reglas de política y multiplicadores}
\lecturas{\BE{9--10} \quad \BE{Apéndice A.3} \quad \FA{10} \quad \FA{13}}
\consultalib{Fair (2024), caps.~9, 11, 12, 14 y 15}
\nota{Canales de transmisión, evaluación de reglas de tasa de interés,
multiplicadores del gasto público y efectos de los choques de inflación.}

\tema{14.6 Pronóstico con modelos macroeconométricos y su evaluación}
\lecturas{\BE{11}}
\nota{Modelos con corrección de errores frente a VAR en diferencias y la
falla de pronóstico ante cambios de media. Es el puente hacia el Bloque 15,
donde la evaluación y la combinación de pronósticos se estudian con detalle.}

%------------------------------------------------------------------------------
\bloque{Bloque 15 --- Pronóstico económico: evaluación, combinación y \emph{nowcasting} --- 4 semanas}
{\textbf{Bloque nuevo.} Reúne los modelos de todo el plan bajo la pregunta del
pronóstico: cómo se diseña un ejercicio fuera de muestra, con qué referentes
se compara, cómo se prueba si un modelo pronostica mejor que otro, cómo se
combinan pronósticos y cómo se evalúan los pronósticos de densidad. Cierra con
el \emph{nowcasting}, la herramienta con la que los bancos centrales leen la
coyuntura. Elliott y Timmermann es la lectura principal y Pesaran (cap.~17) la
profundización. Conviene acompañar cada semana con un ejercicio en datos del
BCRP: inflación y PBI, con la Encuesta de Expectativas Macroeconómicas como
referente.}

\tema{15.1 Funciones de pérdida, pronóstico óptimo y reglas de puntuación}
\lecturas{\ET{2.1--2.5} \quad \ET{3.1--3.4}}
\apoyo{\PS{17.1--17.5}}
\nota{Pérdida cuadrática, absoluta y asimétrica (lin-lin y linex); el
pronóstico óptimo bajo cada una; pronósticos puntuales, de intervalo y de
densidad; reglas de puntuación para distribuciones; incertidumbre de
parámetros en el enfoque clásico y en el bayesiano.}

\tema{15.2 Diseño fuera de muestra, referentes y datos en tiempo real}
\lecturas{\ET{16.1--16.2} \quad \ET{17.8} \quad \ET{21.3}}
\apoyo{\PS{17.7} \quad \PS{17.9}}
\consultalib{Stock y Watson (2012), caps.~14 y 16.2}
\articulosB{Atkeson y Ohanian (2001) \quad Croushore y Stark (2001) \quad
Marcellino, Stock y Watson (2006)}
\nota{Ventanas recursivas y móviles, elección del punto de corte de la
muestra, pronósticos multiperiodo iterados y directos. Referentes que hay que
batir: paseo aleatorio, AR, media móvil de Atkeson y Ohanian y pronósticos de
encuestas. Datos en tiempo real: por qué evaluar con las cifras disponibles en
cada fecha y no con las revisadas (las revisiones del PBI del INEI y del BCRP
son un buen caso).}

\tema{15.3 Selección de modelos y pronóstico univariado}
\lecturas{\ET{6.1--6.5} \quad \ET{7.3--7.5} \quad \ET{20.1--20.2}}
\consultalib{Cowpertwait y Metcalfe (2009), cap.~3 \quad Cryer y Chan (2008),
caps.~9 y 10.5 \quad Hyndman y Athanasopoulos (2021), caps.~5, 8 y 9}
\nota{Criterios de información frente a validación cruzada para series de
tiempo; \emph{lasso} como selector. Pronóstico con ARIMA y SARIMA,
suavizamiento exponencial y Holt--Winters (modelos ETS) y componentes no
observados; pronóstico de variables con tendencia. Retoma los Bloques 2 y 4
desde la óptica del pronóstico.}

\tema{15.4 Pronóstico multivariado y con muchos predictores}
\lecturas{\ET{9.2--9.6} \quad \ET{10.1--10.5} \quad \ET{6.6} \quad \ET{11.3}}
\apoyo{\PS{21.7}}
\consultalib{Hansen (2021), caps.~29.5--29.10 y 29.17 \quad Hastie, Tibshirani
y Friedman (2009), caps.~3.4, 10 y 15 \quad Sosa Escudero (2024) \quad Del
Negro y Schorfheide (2013) \quad Pesaran (2015), cap.~33.6}
\articulosB{Bańbura, Giannone y Reichlin (2010) \quad Medeiros, Vasconcelos,
Veiga y Zilberman (2021)}
\nota{VAR, VAR bayesianos grandes, DSGE y pronóstico condicional; factores y
componentes principales como forma de usar cientos de predictores;
\emph{ridge}, \emph{lasso}, \emph{bagging}, bosques aleatorios y
\emph{boosting} aplicados al pronóstico de la inflación, siempre comparados
con los referentes del tema 15.2. Las técnicas de aprendizaje automático se
estudian en el tema 13.6 del sílabo de econometría; aquí solo se aplican.}

\tema{15.5 Evaluación y comparación de pronósticos}
\lecturas{\ET{15.1--15.3} \quad \ET{16.3--16.4} \quad \ET{17.1--17.7} \quad \ET{17.10}}
\apoyo{\PS{17.10--17.12}}
\articulosB{Diebold y Mariano (1995) \quad Clark y West (2007) \quad
Giacomini y White (2006)}
\consultalib{Mincer y Zarnowitz (1969) \quad West (1996) \quad Hansen, Lunde
y Nason (2011) \quad Diebold (2015)}
\nota{Insesgamiento y eficiencia: regresiones de Mincer--Zarnowitz y pruebas
de racionalidad. Pruebas de \emph{encompassing}. Igualdad de pérdida esperada
con la prueba de Diebold--Mariano; por qué falla con modelos anidados y la
corrección de Clark--West; la habilidad predictiva condicional de
Giacomini--White. Comparación de muchos modelos y minería de datos:
\emph{reality check}, SPA y \emph{model confidence set}. Comparación dentro y
fuera de muestra.}

\tema{15.6 Combinación de pronósticos e inestabilidad}
\lecturas{\ET{14.1--14.4} \quad \ET{14.6--14.7} \quad \ET{19.1--19.3} \quad \ET{19.6--19.7}}
\apoyo{\PS{17.8}}
\articulosB{Bates y Granger (1969) \quad Stock y Watson (2004)}
\consultalib{Timmermann (2006) \quad Rossi (2013) \quad Clements y Hendry
(1998) \quad Stock y Watson (2007)}
\nota{Pesos óptimos y su estimación; por qué el promedio simple suele ganar
(el «puzzle de la combinación»); promedios bayesianos de modelos. Quiebres y
desempeño predictivo: ventanas móviles, correcciones de intercepto y
combinación como defensa frente a la inestabilidad. Enlaza con los temas 4.5
y 12.4 y con la falla de pronóstico del tema 14.6.}

\tema{15.7 Pronósticos de densidad, intervalos y \emph{fan charts}}
\lecturas{\ET{13.4--13.5} \quad \ET{14.5} \quad \ET{18.1--18.3} \quad \ET{18.5}}
\apoyo{\PS{17.14}}
\articulosB{Diebold, Gunther y Tay (1998) \quad Gneiting y Raftery (2007)}
\consultalib{Britton, Fisher y Whitley (1998)}
\nota{Cómo se construye un pronóstico de densidad (simulación, bootstrap,
distribución predictiva bayesiana y regresión cuantílica) y cómo se comunica
en un \emph{fan chart}, como los del Reporte de Inflación del BCRP.
Evaluación con la transformada integral de probabilidad (PIT), el puntaje
logarítmico y el CRPS; evaluación de intervalos y combinación de densidades.
Permite evaluar las densidades del \emph{growth at risk} (10.4) y de los BVAR
(Bloque 12).}

\tema{15.8 \emph{Nowcasting} y datos de frecuencia mixta}
\lecturas{\ET{21.4}}
\articulosB{Giannone, Reichlin y Small (2008) \quad Ghysels, Sinko y Valkanov
(2007) \quad Bańbura, Giannone, Modugno y Reichlin (2013)}
\consultalib{Baffigi, Golinelli y Parigi (2004) \quad Dagum y Cholette (2006),
caps.~6--7}
\nota{Ecuaciones puente, regresiones MIDAS y modelos de factores dinámicos con
el filtro de Kalman y datos faltantes (el borde irregular de los datos);
contenido informativo de cada publicación. Desagregación temporal (Denton,
Chow--Lin) para pasar de datos anuales o trimestrales a mensuales. Aplicación:
\emph{nowcast} del PBI trimestral peruano con el PBI mensual y los indicadores
adelantados del BCRP. Retoma el Bloque 8 y el tema 12.5.}

%------------------------------------------------------------------------------
\newpage
\section*{Anexo A. De los capítulos a los bloques}
\addcontentsline{toc}{section}{Anexo A. De los capítulos a los bloques}

Capítulos de cada libro que entran en el plan y temas donde se leen. Las
secciones excluidas se detallan en el Anexo B.

\newcommand{\mapa}[2]{%
\subsection*{#1}
{\small\filasalternas
\noindent\begin{tabularx}{\textwidth}{>{\raggedright\arraybackslash}p{1.1cm}X>{\raggedright\arraybackslash}p{3.8cm}}
\toprule
\rowcolor{tonoOscuro}\cab{Cap.} & \cab{Título} & \cab{Temas} \\
\midrule
#2
\bottomrule
\end{tabularx}
}}

\mapa{A.1 Martin, Hurn y Harris, \emph{Econometric Modelling with Time Series} (texto base)}{%
1 & The Maximum Likelihood Principle & P.1 \\
2 & Properties of Maximum Likelihood Estimators (2.1--2.5) & P.1 \\
3 & Numerical Estimation Methods & P.2 \\
4 & Hypothesis Testing & P.2 \\
7 & Autocorrelated Regression Models (7.1--7.7) & 3.2 \\
9 & Quasi-Maximum Likelihood Estimation (9.1--9.4) & 9.1 \\
10 & Generalized Method of Moments & 13.3 \\
12 & Estimation by Simulation (12.1--12.5) & 13.3 \\
13 & Linear Time Series Models & 1.1, 1.4, 2.1, 2.2, 2.3, 5.1 \\
14 & Structural Vector Autoregressions & 6.1, 6.2 \\
15 & Latent Factor Models & 8.1, 8.2, 8.3, 8.5 \\
16 & Nonstationary Distribution Theory & 4.1, 4.2 \\
17 & Unit Root Testing & 4.3, 4.4, 4.5 \\
18 & Cointegration & 7.2, 7.3, 7.4 \\
19 & Nonlinearities in Mean (19.1--19.4, 19.6, 19.8) & 10.1, 10.2, 10.3 \\
20 & Nonlinearities in Variance (20.1--20.6) & 9.1, 9.2, 9.3 \\
}

\mapa{A.2 Bårdsen, Eitrheim, Jansen y Nymoen, \emph{The Econometrics of Macroeconomic Modelling} (texto base)}{%
1 & Introduction & 14.1 \\
2 & Methodological issues of large-scale macromodels & 14.1 \\
3 & Inflation in open economies: the main-course model & 14.3 \\
4 & The Phillips curve & 14.2, 14.3 \\
5 & Wage bargaining and price-setting & 14.3 \\
6 & Wage--price dynamics & 14.3 \\
7 & The New Keynesian Phillips curve & 14.4 \\
8 & Money and inflation & 14.4 \\
9 & Transmission channels and model properties & 14.2, 14.5 \\
10 & Evaluation of monetary policy rules & 14.5 \\
11 & Forecasting using econometric models & 14.6 \\
Ap. & The Lucas critique; rational expectations models; interim multipliers & 13.1, 14.2, 14.5 \\
}

\mapa{A.3 Enders, \emph{Applied Econometric Time Series}}{%
1 & Difference Equations & 1.1, 1.2, 1.3 \\
2 & Stationary Time-Series Models & 1.4, 2.1, 2.2, 2.4, 2.5 \\
3 & Modeling Volatility & 9.1, 9.2, 9.3 \\
4 & Models with Trend (sin 4.11) & 4.1, 4.3, 4.4, 4.5, 4.6 \\
5 & Multiequation Time-Series Models & 3.2, 5.1--5.4, 6.1, 6.2, 6.3 \\
6 & Cointegration and Error-Correction Models & 7.1, 7.2, 7.3, 7.4 \\
7 & Nonlinear Models and Breaks & 4.5, 10.1, 10.2, 10.3 \\
}

\mapa{A.4 Hamilton, \emph{Time Series Analysis}}{%
1 & Difference Equations & 1.2 \\
2 & Lag Operators & 1.3 \\
3 & Stationary ARMA Processes & 1.4, 2.1 \\
4 & Forecasting & 2.2, 2.4 \\
5 & Maximum Likelihood Estimation (sin 5.9) & P.2, 2.3 \\
6 & Spectral Analysis & 3.3, 3.4 \\
7 & Asymptotic Distribution Theory & 3.1 \\
8 & Linear Regression Models & 3.2 \\
10 & Covariance-Stationary Vector Processes (10.1--10.4) & 3.3, 5.1 \\
11 & Vector Autoregressions & 5.2, 5.3, 5.4, 6.1, 6.4 \\
12 & Bayesian Analysis & 11.1, 11.3, 12.1 \\
13 & The Kalman Filter & 8.1--8.4 \\
14 & Generalized Method of Moments & 13.3 \\
15--16 & Nonstationary Time Series; Deterministic Time Trends & 4.1 \\
17 & Univariate Processes with Unit Roots (17.1--17.8) & 4.2, 4.3, 4.4 \\
18 & Unit Roots in Multivariate Time Series & 4.2, 7.1 \\
19 & Cointegration & 7.2 \\
20 & Full-Information ML Analysis of Cointegrated Systems & 7.3, 7.4 \\
21 & Time Series Models of Heteroskedasticity & 9.1, 9.2 \\
22 & Modeling Time Series with Changes in Regime & 10.2 \\
}

\mapa{A.5 Kilian y Lütkepohl, \emph{Structural Vector Autoregressive Analysis}}{%
2 & Vector Autoregressive Models & 5.1, 5.2, 5.3 \\
3 & Vector Error Correction Models & 7.3, 7.4 \\
4 & Structural VAR Tools & 5.4, 6.8 \\
5 & Bayesian VAR Analysis & 12.1 \\
6 & VAR Models and other Macroeconometric Models & 13.1, 14.1 \\
7 & A Historical Perspective on Causal Inference & 5.3, 6.1 \\
8 & Identification by Short-Run Restrictions & 6.1, 6.2 \\
9 & Estimation Subject to Short-Run Restrictions & 6.2, 12.2 \\
10 & Identification by Long-Run Restrictions & 6.3 \\
11 & Estimation Subject to Long-Run Restrictions (11.1--11.3) & 6.3 \\
12 & Inference in Models Identified by Short- or Long-Run Restrictions (12.1--12.8) & 6.4, 6.7, 7.5 \\
13 & Identification by Sign Restrictions (sin 13.8) & 6.5, 12.2 \\
14 & Identification by Heteroskedasticity or Non-Gaussianity & 6.6 \\
15 & Identification Based on Extraneous Data & 6.6 \\
16 & Structural VAR Analysis in a Data-Rich Environment (16.1--16.3) & 8.5, 12.5 \\
18 & Nonlinear Structural VAR Models (18.1--18.5) & 10.3, 12.4 \\
}

\mapa{A.6 Blake y Mumtaz, \emph{Applied Bayesian Econometrics for Central Bankers}}{%
1 & Gibbs Sampling for Linear Regression Models & 11.1, 11.2, 11.3 \\
2 & Gibbs Sampling for Vector Autoregressions & 12.1, 12.2 \\
3 & Gibbs Sampling for State Space Models & 8.2, 12.3, 12.4, 12.5 \\
4 & Gibbs Sampling for Markov Switching Models & 12.5 \\
5 & An Introduction to the Metropolis Hastings Algorithm & 11.4, 12.4 \\
6 & Bayesian Estimation of Linear DSGE Models & 13.2 \\
7 & State-Space Models with Time-Varying Parameters & 12.4 \\
}

\mapa{A.7 Fair, \emph{Macroeconomic Modeling: The Cowles Commission Approach}}{%
1 & Why This Book? & 14.1 \\
3 & The Econometrics & 14.1 \\
4 & Specification and Estimation of the US Model & 14.1 \\
10 & Size of Fed's Effect on Output, Unemployment, and Inflation & 14.5 \\
13 & Size of Government Spending Multipliers & 14.5 \\
}

\mapa{A.8 Elliott y Timmermann, \emph{Economic Forecasting} (Bloque 15)}{%
2 & Loss Functions (2.1--2.5) & 15.1 \\
3 & The Parametric Forecasting Problem (3.1--3.4) & 15.1 \\
6 & Model Selection (6.1--6.6) & 15.3, 15.4 \\
7 & Univariate Linear Prediction Models (7.3--7.5) & 15.3 \\
9 & Vector Autoregressions (9.2--9.6) & 15.4 \\
10 & Forecasting in a Data-Rich Environment (10.1--10.5) & 15.4 \\
11 & Nonparametric Forecasting Methods (11.3) & 15.4 \\
13 & Volatility and Density Forecasting (13.4--13.5) & 15.7 \\
14 & Forecast Combinations (sin 14.8) & 15.6, 15.7 \\
15 & Desirable Properties of Forecasts (15.1--15.3) & 15.5 \\
16 & Evaluation of Individual Forecasts (16.1--16.4) & 15.2, 15.5 \\
17 & Evaluation and Comparison of Multiple Forecasts (17.1--17.8, 17.10) & 15.2, 15.5 \\
18 & Evaluating Density Forecasts (18.1--18.3, 18.5) & 15.7 \\
19 & Forecasting under Model Instability (19.1--19.3, 19.6--19.7) & 15.6 \\
20 & Trending Variables and Forecasting (20.1--20.2) & 15.3 \\
21 & Forecasting Nonstandard Data (21.3--21.4) & 15.2, 15.8 \\
}

\mapa{A.9 Pesaran, \emph{Time Series and Panel Data Econometrics} (profundización)}{%
17 & Introduction to Forecasting (17.1--17.5, 17.7--17.12, 17.14) & 15.1, 15.2, 15.5, 15.6, 15.7 \\
21 & Vector Autoregressive Models (21.7) & 15.4 \\
}

%------------------------------------------------------------------------------
\section*{Anexo B. Capítulos y secciones que no entran en el plan}
\addcontentsline{toc}{section}{Anexo B. Capítulos y secciones que no entran en el plan}

{\small\filasalternas
\noindent\begin{tabularx}{\textwidth}{>{\raggedright\arraybackslash}p{4.6cm}X}
\toprule
\rowcolor{tonoOscuro}\cab{Materia} & \cab{Capítulos y secciones} \\
\midrule
Ya cubierto en el sílabo de econometría & \MHH{5--6} y \MHH{8} (regresión lineal, no lineal y heteroscedástica), \MHH{11} (métodos no paramétricos), \MHH{21} (series de tiempo discretas); \HM{9} (ecuaciones simultáneas) \\
Datos de panel & \EN{4.11} (raíces unitarias en panel); panel VAR de 1ECO78. Ver Anexo D.2 \\
Finanzas & Valor en riesgo, pérdida esperada, \emph{backtesting} y CoVaR de 1ECO78. Ver Anexo D.1 \\
Otras secciones & \MHH{7.8}, \MHH{12.6} (modelos en tiempo continuo), \MHH{19.5}, \MHH{19.7}; \HM{5.9}, \HM{10.5}, \HM{17.9}; \KL{11.4--11.5}, \KL{12.9}, \KL{13.8}, \KL{16.4--16.5}, \KL{17} (choques no fundamentales), \KL{18.6--18.9}, \KL{19} (consulta en el tema 4.5) \\
Fair, capítulos de consulta & \FA{2}, \FA{5--9}, \FA{11--12}, \FA{14--16} (citados como consulta); \FA{17--18} (modelos multipaís y material adicional) \\
Elliott y Timmermann & \ET{1} (introducción), \ET{4} (estimación basada en la pérdida), \ET{5} (consulta en el tema 12.1), \ET{8} (los modelos no lineales del Bloque 10 bastan), \ET{11.1--11.2}, \ET{12} (pronósticos binarios), \ET{13.1--13.3} y \ET{13.6--13.7} (volatilidad, propia de finanzas), \ET{18.4}, \ET{20.3--20.5}, \ET{21.1--21.2} (conteos y duraciones) \\
Pesaran & Todo salvo los capítulos 17 y 21.7; el resto del libro repite temas del plan o del sílabo de econometría \\
\bottomrule
\end{tabularx}
}

\medskip
\noindent Hamilton tiene dos copias en la carpeta. La de la carpeta Time Series
es un escaneo sin texto; la que se usa en el plan es la de la carpeta de
econometría (edición de Princeton, versión de 2020). El texto de Favero
(2001) se ve bien en pantalla pero su texto no se puede extraer, por lo que
solo se cita como consulta.

%------------------------------------------------------------------------------
\section*{Anexo C. Programas de la PUCP utilizados}
\addcontentsline{toc}{section}{Anexo C. Programas de la PUCP utilizados}

\subsection*{C.1 Econometría Intermedia: Macro (1ECO24, pregrado; 2025-1, 2025-2 y 2026-1)}
{\small\filasalternas
\noindent\begin{tabularx}{\textwidth}{>{\raggedright\arraybackslash}p{6.2cm}X}
\toprule
\rowcolor{tonoOscuro}\cab{Tema del programa} & \cab{Temas del plan} \\
\midrule
Características de los datos macroeconómicos & 1.1 \\
Ecuaciones en diferencias y operadores de rezagos & 1.2, 1.3 \\
Teoría asintótica para series temporales & P, 3.1 \\
Procesos ARMA, pronósticos y máxima verosimilitud & Bloque 2 \\
Tendencias determinísticas y raíces unitarias & Bloque 4 \\
VAR y SVAR & Bloques 5 y 6 \\
Proyecciones locales (2025-2) & 6.7 \\
Cointegración y VECM & Bloque 7 \\
Espacio de estados, filtro de Kalman y otros filtros & 3.4, Bloque 8 \\
Modelos de factores y FAVAR (2025-2) & 8.5 \\
Modelos de volatilidad & Bloque 9 \\
Umbrales y cambio de régimen (opcionales) & 10.1, 10.2 \\
Inteligencia artificial y aprendizaje automático (2025-2) & 10.3 y tema 13.6 del sílabo de econometría \\
\bottomrule
\end{tabularx}
}

\subsection*{C.2 Econometría Avanzada: Series de Tiempo (1ECO78, maestría; 2024-2, 2025-2 y 2026-2)}
{\small\filasalternas
\noindent\begin{tabularx}{\textwidth}{>{\raggedright\arraybackslash}p{6.2cm}X}
\toprule
\rowcolor{tonoOscuro}\cab{Módulo o sesión} & \cab{Temas del plan} \\
\midrule
M1. ARMA, VAR, SVAR e introducción al filtro de Kalman & Bloques 2, 5 y 6; 8.1--8.3 \\
M2. Asintótica, raíces unitarias y cointegración & 3.1, Bloques 4 y 7 \\
M2. Análisis espectral y modelos de factores (2024-2) & 3.3, 8.5 \\
M2. Medidas de riesgo (2025-2, 2026-2) & Anexo D.1 \\
M2. Regresión cuantílica, CoVaR y \emph{growth at risk} & 10.4; Anexo D.1 \\
M3. Econometría bayesiana y BVAR & Bloques 11 y 12 \\
M3. Proyecciones locales & 6.7 \\
M3. Espacio de estados bayesiano, TVP-VAR-SV y \emph{nowcasting} & 12.3, 12.4, 12.5 \\
M3. Factores dinámicos y FAVAR & 8.5, 12.5 \\
M3. Panel VAR & Anexo D.2 \\
M3. Estimación de DSGE y filtros no lineales & Bloque 13 \\
M3. Aprendizaje automático & 10.3 \\
M4. MCMC, volatilidad bayesiana y volatilidad estocástica & 11.3, 11.4, 12.4 \\
M4. Carter--Kohn y Durbin--Koopman & 12.3 \\
\bottomrule
\end{tabularx}
}

\subsection*{C.3 Otros programas}
\begin{itemize}[leftmargin=1.4em,itemsep=1pt]
  \item \textbf{Econometría Intermedia (2ECO01, maestría, 2026-1), Unidad 4}:
        ARMA (Bloque 2), VAR y SVAR (Bloques 5 y 6), raíces unitarias y
        cointegración (Bloques 4 y 7). Las Unidades 1 a 3 están en el sílabo
        de econometría.
  \item \textbf{Macroeconometría (BCRP, 2004)}, 51.º Curso de Extensión
        Universitaria, E.~Morón: VAR estructural (Bloque 6), GMM y calibración
        (13.3), espacio de estados (Bloque 8), regresiones de umbral y de
        transición suave (10.1).
\end{itemize}

%------------------------------------------------------------------------------
\section*{Anexo D. Material opcional}
\addcontentsline{toc}{section}{Anexo D. Material opcional}

Todo lo de este anexo es opcional: no entra en el cómputo de horas ni modifica
el cronograma.

\subsection*{D.1 Medidas de riesgo financiero}
Sigue las sesiones 2.3 y 2.4 de 1ECO78. Momento sugerido: después de los
Bloques 9 y 10.
\begin{itemize}[leftmargin=1.4em,itemsep=1pt]
  \item Valor en riesgo, pérdida esperada y \emph{backtesting}: Tsay (2010),
        cap.~7; Christoffersen (2009).
  \item Riesgo sistémico y CoVaR: Adrian y Brunnermeier (2016); Acharya,
        Pedersen, Philippon y Richardson (2010); Brownlees y Engle (2012).
\end{itemize}

\subsection*{D.2 Panel VAR y raíces unitarias en panel}
Sigue la sesión 3.3 de 1ECO78 y requiere la carpeta Panel Data. Punto de
partida: \EN{4.11}.

\subsection*{D.3 Software}
\begin{itemize}[leftmargin=1.4em,itemsep=1pt]
  \item \textbf{EViews y R}: bases de datos y archivos de Enders (carpeta
        Enders); Pfaff (2008) para raíces unitarias y cointegración en R;
        Shumway y Stoffer para series de tiempo en R.
  \item \textbf{Matlab}: códigos de Blake y Mumtaz y su apéndice de
        introducción a Matlab (\BM{8}).
  \item \textbf{JMulTi}: Lütkepohl y Krätzig (2004).
  \item \textbf{Pronóstico (Bloque 15)}: en R, \texttt{forecast} o
        \texttt{fable} (ARIMA y ETS), \texttt{midasr} (regresiones MIDAS),
        \texttt{scoringRules} (puntaje logarítmico y CRPS) y \texttt{MCS}
        (\emph{model confidence set}); en Python, \texttt{statsmodels}.
  \item \textbf{Datos}: el complemento BCRPData para descargar series del BCRP;
        la Encuesta de Expectativas Macroeconómicas del BCRP como referente de
        pronóstico.
\end{itemize}

\subsection*{D.4 Proyecto aplicado opcional: política monetaria y fiscal en el Perú}
Consolida los Bloques 4 a 7, 12 y 15. Momento sugerido: los pasos 1 a 5
después del Bloque 12 y el paso 6 después del Bloque 15; no tiene horas
asignadas.
\begin{enumerate}[leftmargin=1.6em,itemsep=2pt]
  \item \textbf{Datos.} Series mensuales o trimestrales del BCRP: actividad,
        inflación, tasa de interés de referencia, tipo de cambio y crédito.
  \item \textbf{Propiedades univariadas.} Raíces unitarias y quiebres (Bloque 4).
  \item \textbf{VAR estructural.} Identificación recursiva y por signos; comparar
        con proyecciones locales (Bloque 6).
  \item \textbf{Largo plazo.} Cointegración entre las variables nominales
        (Bloque 7).
  \item \textbf{BVAR.} Estimación bayesiana y pronóstico condicional
        (Bloque 12).
  \item \textbf{Pronóstico.} Ejercicio pseudo fuera de muestra de la
        inflación y del PBI, de preferencia con datos en tiempo real; prueba
        de Diebold--Mariano (o de Clark--West) frente al paseo aleatorio y a
        la Encuesta de Expectativas del BCRP; combinación de los pronósticos
        del ARIMA, el BVAR y el modelo de factores; \emph{fan chart} evaluado
        con la PIT; y \emph{nowcast} del PBI trimestral con el PBI mensual
        (Bloque 15).
  \item \textbf{Referencias peruanas} (1ECO24, 2026-1): Lahura (2017) sobre el
        traspaso de la tasa de referencia; Castillo y Lahura (2022) y Aguilar
        y Lahura (2026) sobre política fiscal.
\end{enumerate}

%------------------------------------------------------------------------------
\section*{Anexo E. Referencias}
\addcontentsline{toc}{section}{Anexo E. Referencias}

\subsection*{E.1 Textos del plan}
{\small
\begin{itemize}[leftmargin=1.4em,itemsep=1pt,label={}]
\item Bårdsen, G., Ø. Eitrheim, E. S. Jansen y R. Nymoen (2005). \emph{The Econometrics of Macroeconomic Modelling}. Oxford University Press.
\item Blake, A. y H. Mumtaz (2017). \emph{Applied Bayesian Econometrics for Central Bankers}, edición actualizada. Centre for Central Banking Studies, Bank of England.
\item Elliott, G. y A. Timmermann (2016). \emph{Economic Forecasting}. Princeton University Press.
\item Enders, W. (2014). \emph{Applied Econometric Time Series}, 4.ª ed. Wiley.
\item Fair, R. C. (2024). \emph{Macroeconomic Modeling: The Cowles Commission Approach}. MIT Press.
\item Hamilton, J. D. (1994). \emph{Time Series Analysis}. Princeton University Press.
\item Kilian, L. y H. Lütkepohl (2017). \emph{Structural Vector Autoregressive Analysis}. Cambridge University Press.
\item Martin, V., S. Hurn y D. Harris (2013). \emph{Econometric Modelling with Time Series: Specification, Estimation and Testing}. Cambridge University Press.
\item Montiel Olea, J. L., M. Plagborg-Møller, E. Qian y C. K. Wolf (2025). Local Projections or VARs? A Primer for Macroeconomists. Documento de trabajo, versión del 21 de marzo de 2025.
\item Pesaran, M. H. (2015). \emph{Time Series and Panel Data Econometrics}. Oxford University Press (carpeta Panel Data).
\end{itemize}
}

\subsection*{E.2 Libros de consulta de la carpeta (sin horas)}
{\small
\begin{itemize}[leftmargin=1.4em,itemsep=1pt,label={}]
\item Brockwell, P. J. y R. A. Davis (2009). \emph{Time Series: Theory and Methods}, 2.ª ed. Springer.
\item Chen, R. y R. S. Tsay (2019). \emph{Nonlinear Time Series Analysis}. Wiley.
\item Cowpertwait, P. S. P. y A. V. Metcalfe (2009). \emph{Introductory Time Series with R}. Springer.
\item Cryer, J. D. y K.-S. Chan (2008). \emph{Time Series Analysis: With Applications in R}, 2.ª ed. Springer (carpeta de econometría).
\item Dagum, E. B. y P. A. Cholette (2006). \emph{Benchmarking, Temporal Distribution, and Reconciliation Methods for Time Series}. Lecture Notes in Statistics 186. Springer (carpeta de econometría).
\item Davidson, R. y J. G. MacKinnon (2004). \emph{Econometric Theory and Methods}. Oxford University Press (carpeta de econometría).
\item Favero, C. A. (2001). \emph{Applied Macroeconometrics}. Oxford University Press.
\item Garratt, A., K. Lee, M. H. Pesaran e Y. Shin (2006). \emph{Global and National Macroeconometric Modelling: A Long-Run Structural Approach}. Oxford University Press.
\item Greene, W. H. (2012). \emph{Econometric Analysis}, 7.ª ed. Pearson (carpeta de econometría).
\item Hansen, B. E. (2021). \emph{Econometrics}. Manuscrito, University of Wisconsin (carpeta de econometría).
\item Hastie, T., R. Tibshirani y J. Friedman (2009). \emph{The Elements of Statistical Learning}, 2.ª ed. Springer (carpeta de econometría).
\item Hayashi, F. (2000). \emph{Econometrics}. Princeton University Press (carpeta de econometría).
\item Hyndman, R. J. y G. Athanasopoulos (2021). \emph{Forecasting: Principles and Practice}, 3.ª ed. OTexts (de libre acceso en línea; no está en la carpeta).
\item Lütkepohl, H. (2005). \emph{New Introduction to Multiple Time Series Analysis}. Springer.
\item Lütkepohl, H. y M. Krätzig (2004). \emph{Applied Time Series Econometrics}. Cambridge University Press.
\item McElreath, R. (2016). \emph{Statistical Rethinking: A Bayesian Course with Examples in R and Stan}. CRC Press (carpeta de econometría).
\item Neusser, K. (2016). \emph{Time Series Econometrics}. Springer.
\item Pfaff, B. (2008). \emph{Analysis of Integrated and Cointegrated Time Series with R}, 2.ª ed. Springer.
\item Shumway, R. H. y D. S. Stoffer (2017). \emph{Time Series Analysis and Its Applications: With R Examples}, 4.ª ed. Springer.
\item Sosa Escudero, W. (2024). \emph{Machine learning para economistas: ideas básicas}. Notas de clase, UNLP (carpeta de econometría).
\item Stock, J. H. y M. W. Watson (2012). \emph{Introducción a la econometría}, 3.ª ed. Pearson (carpeta de econometría).
\item Tsay, R. S. (2010). \emph{Analysis of Financial Time Series}, 3.ª ed. Wiley.
\item Tsay, R. S. (2013). \emph{Multivariate Time Series Analysis: With R and Financial Applications}. Wiley.
\item White, H. (2001). \emph{Asymptotic Theory for Econometricians}, ed. revisada. Academic Press (carpeta de econometría).
\end{itemize}
}

\subsection*{E.3 Referencias tomadas de los programas de la PUCP y del BCRP}
{\small
\begin{itemize}[leftmargin=1.4em,itemsep=1pt,label={}]
\item Acharya, V., L. Pedersen, T. Philippon y M. Richardson (2010). Measuring Systemic Risk. Documento de trabajo.
\item Adrian, T. y M. K. Brunnermeier (2016). CoVaR. \emph{American Economic Review}, 106(7), 1705--1741.
\item Adrian, T., N. Boyarchenko y D. Giannone (2019). Vulnerable Growth. \emph{American Economic Review}, 109(4), 1263--1289.
\item Aguilar, J. y E. Lahura (2026). When Expectations Matter: The Role of Fiscal Foresight in Government Spending Shocks in Peru. \emph{Economic Analysis and Policy}, 89, 505--526.
\item Bernanke, B. y I. Mihov (1998). Measuring Monetary Policy. \emph{Quarterly Journal of Economics}, 113(3), 869--902.
\item Blanchard, O. J. y D. Quah (1989). The Dynamic Effects of Aggregate Demand and Supply Disturbances. \emph{American Economic Review}, 79(4), 655--673.
\item Brownlees, C. y R. Engle (2012). Volatility, Correlation and Tails for Systemic Risk Measurement. Documento de trabajo.
\item Carter, C. K. y R. Kohn (1994). On Gibbs Sampling for State Space Models. \emph{Biometrika}, 81, 541--553.
\item Casella, G. y E. I. George (1992). Explaining the Gibbs Sampler. \emph{The American Statistician}, 46(3), 167--174.
\item Castillo, T. y E. Lahura (2022). Los efectos de la política fiscal sobre la actividad económica en el Perú. \emph{Revista Estudios Económicos}, BCRP, 40, 9--34.
\item Chib, S. (2001). Markov Chain Monte Carlo Methods: Computation and Inference. En J. J. Heckman y E. Leamer (eds.), \emph{Handbook of Econometrics}, vol.~5, cap.~57.
\item Christiano, L. J. y T. J. Fitzgerald (2003). The Band Pass Filter. \emph{International Economic Review}, 44(2), 435--465.
\item Christoffersen, P. F. (2009). Value-at-Risk Models.
\item Cogley, T. y J. M. Nason (1995). Effects of the Hodrick--Prescott Filter on Trend and Difference Stationary Time Series. \emph{Journal of Economic Dynamics and Control}, 19(1--2), 253--278.
\item Cogley, T. y T. J. Sargent (2005). Drifts and Volatilities: Monetary Policies and Outcomes in the Post WWII US. \emph{Review of Economic Dynamics}, 8, 262--302.
\item Durbin, J. y S. J. Koopman (2002). A Simple and Efficient Simulation Smoother for State Space Time Series Analysis. \emph{Biometrika}, 89(3), 603--615.
\item Elliott, G., T. J. Rothenberg y J. H. Stock (1996). Efficient Tests for an Autoregressive Unit Root. \emph{Econometrica}, 64(4), 813--836.
\item Engle, R. F. (2002). Dynamic Conditional Correlation: A Simple Class of Multivariate GARCH Models. \emph{Journal of Business \& Economic Statistics}, 20(3), 339--350.
\item Fernández-Villaverde, J. y J. Rubio-Ramírez (2007). Estimating Macroeconomic Models: A Likelihood Approach. \emph{Review of Economic Studies}, 74, 1059--1087.
\item Gelman, A., J. Carlin, H. Stern, D. Dunson, A. Vehtari y D. Rubin (2013). \emph{Bayesian Data Analysis}, 3.ª ed. CRC Press.
\item Harvey, A. C. (1990). \emph{Forecasting, Structural Time Series Models and the Kalman Filter}. Cambridge University Press.
\item Herbst, E. y F. Schorfheide (2016). \emph{Bayesian Estimation of DSGE Models}. Princeton University Press.
\item Jordà, Ò. (2005). Estimation and Inference of Impulse Responses by Local Projections. \emph{American Economic Review}, 95(1), 161--182.
\item Kim, C.-J. y C. R. Nelson (1999). \emph{State-Space Models with Regime Switching}. MIT Press.
\item Kim, S., N. Shephard y S. Chib (1998). Stochastic Volatility: Likelihood Inference and Comparison with ARCH Models. \emph{Review of Economic Studies}, 65(3), 361--393.
\item Koop, G. (2003). \emph{Bayesian Econometrics}. Wiley.
\item Koop, G. y D. Korobilis (2010). Bayesian Multivariate Time Series Methods for Empirical Macroeconomics. \emph{Foundations and Trends in Econometrics}, 3(4), 267--358.
\item Lahura, E. (2017). El efecto traspaso de la tasa de interés de política monetaria. \emph{Revista Estudios Económicos}, BCRP, 33, 9--27.
\item Litterman, R. (1986). Forecasting with Bayesian Vector Autoregressions: Five Years of Experience. \emph{Journal of Business \& Economic Statistics}, 4(1), 25--38.
\item Sims, C. A. (1980). Macroeconomics and Reality. \emph{Econometrica}, 48(1), 1--48.
\item Stock, J. H. y M. W. Watson (2001). Vector Autoregressions. \emph{Journal of Economic Perspectives}, 15(4), 101--115.
\item Stock, J. H. y M. W. Watson (2002). Forecasting Using Principal Components from a Large Number of Predictors. \emph{Journal of the American Statistical Association}, 97, 1167--1179.
\item Uhlig, H. (2005). What Are the Effects of Monetary Policy on Output? Results from an Agnostic Identification Procedure. \emph{Journal of Monetary Economics}, 52, 381--419.
\end{itemize}
}

\subsection*{E.4 Referencias del Bloque 15 (pronóstico)}
{\small
\begin{itemize}[leftmargin=1.4em,itemsep=1pt,label={}]
\item Atkeson, A. y L. E. Ohanian (2001). Are Phillips Curves Useful for Forecasting Inflation? \emph{Federal Reserve Bank of Minneapolis Quarterly Review}, 25(1), 2--11.
\item Baffigi, A., R. Golinelli y G. Parigi (2004). Bridge Models to Forecast the Euro Area GDP. \emph{International Journal of Forecasting}, 20(3), 447--460.
\item Bańbura, M., D. Giannone, M. Modugno y L. Reichlin (2013). Now-Casting and the Real-Time Data Flow. En G. Elliott y A. Timmermann (eds.), \emph{Handbook of Economic Forecasting}, vol.~2A, cap.~4. Elsevier.
\item Bańbura, M., D. Giannone y L. Reichlin (2010). Large Bayesian Vector Auto Regressions. \emph{Journal of Applied Econometrics}, 25(1), 71--92.
\item Bates, J. M. y C. W. J. Granger (1969). The Combination of Forecasts. \emph{Operational Research Quarterly}, 20(4), 451--468.
\item Britton, E., P. Fisher y J. Whitley (1998). The Inflation Report Projections: Understanding the Fan Chart. \emph{Bank of England Quarterly Bulletin}, 38(1), 30--37.
\item Clark, T. E. y K. D. West (2007). Approximately Normal Tests for Equal Predictive Accuracy in Nested Models. \emph{Journal of Econometrics}, 138(1), 291--311.
\item Clements, M. P. y D. F. Hendry (1998). \emph{Forecasting Economic Time Series}. Cambridge University Press.
\item Croushore, D. y T. Stark (2001). A Real-Time Data Set for Macroeconomists. \emph{Journal of Econometrics}, 105(1), 111--130.
\item Del Negro, M. y F. Schorfheide (2013). DSGE Model-Based Forecasting. En G. Elliott y A. Timmermann (eds.), \emph{Handbook of Economic Forecasting}, vol.~2A, cap.~2. Elsevier.
\item Diebold, F. X. (2015). Comparing Predictive Accuracy, Twenty Years Later: A Personal Perspective on the Use and Abuse of Diebold--Mariano Tests. \emph{Journal of Business \& Economic Statistics}, 33(1), 1--9.
\item Diebold, F. X., T. A. Gunther y A. S. Tay (1998). Evaluating Density Forecasts with Applications to Financial Risk Management. \emph{International Economic Review}, 39(4), 863--883.
\item Diebold, F. X. y R. S. Mariano (1995). Comparing Predictive Accuracy. \emph{Journal of Business \& Economic Statistics}, 13(3), 253--263.
\item Ghysels, E., A. Sinko y R. Valkanov (2007). MIDAS Regressions: Further Results and New Directions. \emph{Econometric Reviews}, 26(1), 53--90.
\item Giacomini, R. y H. White (2006). Tests of Conditional Predictive Ability. \emph{Econometrica}, 74(6), 1545--1578.
\item Giannone, D., L. Reichlin y D. Small (2008). Nowcasting: The Real-Time Informational Content of Macroeconomic Data. \emph{Journal of Monetary Economics}, 55(4), 665--676.
\item Gneiting, T. y A. E. Raftery (2007). Strictly Proper Scoring Rules, Prediction, and Estimation. \emph{Journal of the American Statistical Association}, 102(477), 359--378.
\item Hansen, P. R., A. Lunde y J. M. Nason (2011). The Model Confidence Set. \emph{Econometrica}, 79(2), 453--497.
\item Marcellino, M., J. H. Stock y M. W. Watson (2006). A Comparison of Direct and Iterated Multistep AR Methods for Forecasting Macroeconomic Time Series. \emph{Journal of Econometrics}, 135(1--2), 499--526.
\item Medeiros, M. C., G. F. R. Vasconcelos, Á. Veiga y E. Zilberman (2021). Forecasting Inflation in a Data-Rich Environment: The Benefits of Machine Learning Methods. \emph{Journal of Business \& Economic Statistics}, 39(1), 98--119.
\item Mincer, J. y V. Zarnowitz (1969). The Evaluation of Economic Forecasts. En J. Mincer (ed.), \emph{Economic Forecasts and Expectations}. NBER.
\item Rossi, B. (2013). Advances in Forecasting under Instability. En G. Elliott y A. Timmermann (eds.), \emph{Handbook of Economic Forecasting}, vol.~2B, cap.~21. Elsevier.
\item Stock, J. H. y M. W. Watson (2004). Combination Forecasts of Output Growth in a Seven-Country Data Set. \emph{Journal of Forecasting}, 23(6), 405--430.
\item Stock, J. H. y M. W. Watson (2007). Why Has U.S. Inflation Become Harder to Forecast? \emph{Journal of Money, Credit and Banking}, 39(s1), 3--33.
\item Timmermann, A. (2006). Forecast Combinations. En G. Elliott, C. W. J. Granger y A. Timmermann (eds.), \emph{Handbook of Economic Forecasting}, vol.~1, cap.~4. Elsevier.
\item West, K. D. (1996). Asymptotic Inference about Predictive Ability. \emph{Econometrica}, 64(5), 1067--1084.
\end{itemize}
}

\noindent\small Los artículos de esta sección no están en la carpeta; los
\emph{Handbooks} y las revistas suelen estar disponibles por la biblioteca
universitaria o en versiones de documento de trabajo de sus autores.
\normalsize

\vspace{1.2cm}
\franjas[2pt]
\begin{center}\small\color{tonoPizarra}
Plan de lectura de elaboración propia. Los títulos de los capítulos
corresponden a las obras citadas y se reproducen únicamente con fines de
referencia bibliográfica. La asignación de capítulos y secciones se basa en
los índices de los libros de la carpeta; los temas de la PUCP proceden de los
programas citados en el Anexo C.
\end{center}

\end{document}
